% Môn: Vật lý
% Lớp: 10
% Chương: Chương III. Động lực học
% Bài: L10C3 Bài 13. Tổng hợp và phân tích lực. Cân bằng lực
% Số câu: 154
% App đọc file này trực tiếp — sửa rồi Commit trên GitHub, bấm Đồng bộ GitHub .tex

% L10C3 Bài 13. Tổng hợp và phân tích lực. Cân bằng lực

% ===== Câu 1 =====
% ID: AIQ_L10C3_FULL_0009
% Mức: VD
\dangbt{Nhận biết lực, biểu diễn vectơ lực và các đặc trưng của lực}
\begin{ex}
% ID: AIQ_L10C3_FULL_0009
% Mức: VD
Hai lực vuông góc đồng quy có độ lớn lần lượt 3 $N$ và $6\,\text{N}$. Tính độ lớn hợp lực (kết quả theo đơn vị N).
\choice
{\True 6.71}
{7.71}
{5.71}
{13.42}
\loigiai{
Áp dụng hệ thức phù hợp của bài học, thay số được $6.71\,\text{N}$.
}
\end{ex}

% ===== Câu 2 =====
% ID: AIQ_L10C3_FULL_0010
% Mức: VD
\begin{ex}
% ID: AIQ_L10C3_FULL_0010
% Mức: VD
Hai lực vuông góc đồng quy có độ lớn lần lượt 4 $N$ và $4\,\text{N}$. Tính độ lớn hợp lực (kết quả theo đơn vị N).
\choice
{6.66}
{\True 5.66}
{4.66}
{11.32}
\loigiai{
Áp dụng hệ thức phù hợp của bài học, thay số được $5.66\,\text{N}$.
}
\end{ex}

% ===== Câu 3 =====
% ID: AIQ_L10C3_FULL_0011
% Mức: TH
\begin{ex}
% ID: AIQ_L10C3_FULL_0011
% Mức: TH
Hai lực vuông góc đồng quy có độ lớn lần lượt 5 $N$ và $5\,\text{N}$. Tính độ lớn hợp lực (kết quả theo đơn vị N).
\choiceTF
{\True Giá trị cần tìm là $7.07\,\text{N}$.}
{Giá trị cần tìm là $8.07\,\text{N}$.}
{\True Phải xác định đầy đủ các lực hoặc đại lượng liên quan trước khi tính.}
{Có thể bỏ qua điều kiện áp dụng của công thức.}
\loigiai{
\begin{itemchoice}
\item (Đúng) Giá trị cần tìm là $7.07\,\text{N}$. (Vì): kết quả $7.07\,\text{N}$. b) Sai. c) Đúng. d) Sai vì phải kiểm tra điều kiện áp dụng
\item (Sai) Giá trị cần tìm là $8.07\,\text{N}$.
\item (Đúng) Phải xác định đầy đủ các lực hoặc đại lượng liên quan trước khi tính.
\item (Sai) Có thể bỏ qua điều kiện áp dụng của công thức.
\end{itemchoice}
}
\end{ex}

% ===== Câu 4 =====
% ID: AIQ_L10C3_FULL_0012
% Mức: TH
\begin{ex}
% ID: AIQ_L10C3_FULL_0012
% Mức: TH
Hai lực vuông góc đồng quy có độ lớn lần lượt 6 $N$ và $6\,\text{N}$. Tính độ lớn hợp lực (kết quả theo đơn vị N).
\choiceTF
{\True Giá trị cần tìm là $8.49\,\text{N}$.}
{Giá trị cần tìm là $9.49\,\text{N}$.}
{\True Phải xác định đầy đủ các lực hoặc đại lượng liên quan trước khi tính.}
{Có thể bỏ qua điều kiện áp dụng của công thức.}
\loigiai{
\begin{itemchoice}
\item (Đúng) Giá trị cần tìm là $8.49\,\text{N}$. (Vì): kết quả $8.49\,\text{N}$. b) Sai. c) Đúng. d) Sai vì phải kiểm tra điều kiện áp dụng
\item (Sai) Giá trị cần tìm là $9.49\,\text{N}$.
\item (Đúng) Phải xác định đầy đủ các lực hoặc đại lượng liên quan trước khi tính.
\item (Sai) Có thể bỏ qua điều kiện áp dụng của công thức.
\end{itemchoice}
}
\end{ex}

% ===== Câu 5 =====
% ID: AIQ_L10C3_FULL_0013
% Mức: VD
\begin{ex}
% ID: AIQ_L10C3_FULL_0013
% Mức: VD
Hai lực vuông góc đồng quy có độ lớn lần lượt 3 $N$ và $4\,\text{N}$. Tính độ lớn hợp lực (kết quả theo đơn vị N).
\choiceTF
{\True Giá trị cần tìm là $5\,\text{N}$.}
{Giá trị cần tìm là $6\,\text{N}$.}
{\True Phải xác định đầy đủ các lực hoặc đại lượng liên quan trước khi tính.}
{Có thể bỏ qua điều kiện áp dụng của công thức.}
\loigiai{
\begin{itemchoice}
\item (Đúng) Giá trị cần tìm là $5\,\text{N}$. (Vì): kết quả $5\,\text{N}$. b) Sai. c) Đúng. d) Sai vì phải kiểm tra điều kiện áp dụng
\item (Sai) Giá trị cần tìm là $6\,\text{N}$.
\item (Đúng) Phải xác định đầy đủ các lực hoặc đại lượng liên quan trước khi tính.
\item (Sai) Có thể bỏ qua điều kiện áp dụng của công thức.
\end{itemchoice}
}
\end{ex}

% ===== Câu 6 =====
% ID: AIQ_L10C3_FULL_0014
% Mức: VD
\begin{ex}
% ID: AIQ_L10C3_FULL_0014
% Mức: VD
Hai lực vuông góc đồng quy có độ lớn lần lượt 4 $N$ và $5\,\text{N}$. Tính độ lớn hợp lực (kết quả theo đơn vị N).
\choiceTF
{\True Giá trị cần tìm là $6.4\,\text{N}$.}
{Giá trị cần tìm là $7.4\,\text{N}$.}
{\True Phải xác định đầy đủ các lực hoặc đại lượng liên quan trước khi tính.}
{Có thể bỏ qua điều kiện áp dụng của công thức.}
\loigiai{
\begin{itemchoice}
\item (Đúng) Giá trị cần tìm là $6.4\,\text{N}$. (Vì): kết quả $6.4\,\text{N}$. b) Sai. c) Đúng. d) Sai vì phải kiểm tra điều kiện áp dụng
\item (Sai) Giá trị cần tìm là $7.4\,\text{N}$.
\item (Đúng) Phải xác định đầy đủ các lực hoặc đại lượng liên quan trước khi tính.
\item (Sai) Có thể bỏ qua điều kiện áp dụng của công thức.
\end{itemchoice}
}
\end{ex}

% ===== Câu 7 =====
% ID: AIQ_L10C3_FULL_0015
% Mức: VD
\begin{ex}
% ID: AIQ_L10C3_FULL_0015
% Mức: VD
Hai lực vuông góc đồng quy có độ lớn lần lượt 5 $N$ và $6\,\text{N}$. Tính độ lớn hợp lực (kết quả theo đơn vị N).
\choiceTF
{\True Giá trị cần tìm là $7.81\,\text{N}$.}
{Giá trị cần tìm là $8.81\,\text{N}$.}
{\True Phải xác định đầy đủ các lực hoặc đại lượng liên quan trước khi tính.}
{Có thể bỏ qua điều kiện áp dụng của công thức.}
\loigiai{
\begin{itemchoice}
\item (Đúng) Giá trị cần tìm là $7.81\,\text{N}$. (Vì): kết quả $7.81\,\text{N}$. b) Sai. c) Đúng. d) Sai vì phải kiểm tra điều kiện áp dụng
\item (Sai) Giá trị cần tìm là $8.81\,\text{N}$.
\item (Đúng) Phải xác định đầy đủ các lực hoặc đại lượng liên quan trước khi tính.
\item (Sai) Có thể bỏ qua điều kiện áp dụng của công thức.
\end{itemchoice}
}
\end{ex}

% ===== Câu 8 =====
% ID: AIQ_L10C3_FULL_0016
% Mức: TH
\begin{ex}
% ID: AIQ_L10C3_FULL_0016
% Mức: TH
Hai lực vuông góc đồng quy có độ lớn lần lượt 6 $N$ và $4\,\text{N}$. Tính độ lớn hợp lực (kết quả theo đơn vị N).
\shortans{7.21}
\loigiai{
Thay số vào hệ thức của bài học thu được $7.21\,\text{N}$.
}
\end{ex}

% ===== Câu 9 =====
% ID: AIQ_L10C3_FULL_0017
% Mức: TH
\begin{ex}
% ID: AIQ_L10C3_FULL_0017
% Mức: TH
Hai lực vuông góc đồng quy có độ lớn lần lượt 3 $N$ và $5\,\text{N}$. Tính độ lớn hợp lực (kết quả theo đơn vị N).
\shortans{5.83}
\loigiai{
Thay số vào hệ thức của bài học thu được $5.83\,\text{N}$.
}
\end{ex}

% ===== Câu 10 =====
% ID: AIQ_L10C3_FULL_0018
% Mức: TH
\begin{ex}
% ID: AIQ_L10C3_FULL_0018
% Mức: TH
Hai lực vuông góc đồng quy có độ lớn lần lượt 4 $N$ và $6\,\text{N}$. Tính độ lớn hợp lực (kết quả theo đơn vị N).
\shortans{7.21}
\loigiai{
Thay số vào hệ thức của bài học thu được $7.21\,\text{N}$.
}
\end{ex}

% ===== Câu 11 =====
% ID: AIQ_L10C3_FULL_0019
% Mức: VD
\begin{ex}
% ID: AIQ_L10C3_FULL_0019
% Mức: VD
Hai lực vuông góc đồng quy có độ lớn lần lượt 5 $N$ và $4\,\text{N}$. Tính độ lớn hợp lực (kết quả theo đơn vị N).
\shortans{6.4}
\loigiai{
Thay số vào hệ thức của bài học thu được $6.4\,\text{N}$.
}
\end{ex}

% ===== Câu 12 =====
% ID: AIQ_L10C3_FULL_0020
% Mức: VD
\begin{ex}
% ID: AIQ_L10C3_FULL_0020
% Mức: VD
Hai lực vuông góc đồng quy có độ lớn lần lượt 6 $N$ và $5\,\text{N}$. Tính độ lớn hợp lực (kết quả theo đơn vị N).
\shortans{7.81}
\loigiai{
Thay số vào hệ thức của bài học thu được $7.81\,\text{N}$.
}
\end{ex}

% ===== Câu 13 =====
% ID: AIQ_L10C3_FULL_0021
% Mức: VD
\begin{ex}
% ID: AIQ_L10C3_FULL_0021
% Mức: VD
Hai lực vuông góc đồng quy có độ lớn lần lượt 3 $N$ và $6\,\text{N}$. Tính độ lớn hợp lực (kết quả theo đơn vị N).
\shortans{6.71}
\loigiai{
Thay số vào hệ thức của bài học thu được $6.71\,\text{N}$.
}
\end{ex}

% ===== Câu 14 =====
% ID: AIQ_L10C3_FULL_0022
% Mức: VD
\begin{bt}
% ID: AIQ_L10C3_FULL_0022
% Mức: VD
Hai lực vuông góc đồng quy có độ lớn lần lượt 4 $N$ và $4\,\text{N}$. Tính độ lớn hợp lực (kết quả theo đơn vị N).
\loigiai{
Xác định các đại lượng đã cho, chọn hệ thức phù hợp và thay số. Kết quả: $5.66\,\text{N}$.
}
\end{bt}

% ===== Câu 15 =====
% ID: AIQ_L10C3_FULL_0023
% Mức: VD
\begin{bt}
% ID: AIQ_L10C3_FULL_0023
% Mức: VD
Hai lực vuông góc đồng quy có độ lớn lần lượt 5 $N$ và $5\,\text{N}$. Tính độ lớn hợp lực (kết quả theo đơn vị N).
\loigiai{
Xác định các đại lượng đã cho, chọn hệ thức phù hợp và thay số. Kết quả: $7.07\,\text{N}$.
}
\end{bt}

% ===== Câu 16 =====
% ID: AIQ_L10C3_FULL_0024
% Mức: VD
\begin{bt}
% ID: AIQ_L10C3_FULL_0024
% Mức: VD
Hai lực vuông góc đồng quy có độ lớn lần lượt 6 $N$ và $6\,\text{N}$. Tính độ lớn hợp lực (kết quả theo đơn vị N).
\loigiai{
Xác định các đại lượng đã cho, chọn hệ thức phù hợp và thay số. Kết quả: $8.49\,\text{N}$.
}
\end{bt}

% ===== Câu 17 =====
% ID: AIQ_L10C3_FULL_0025
% Mức: VDC
\begin{bt}
% ID: AIQ_L10C3_FULL_0025
% Mức: VDC
Hai lực vuông góc đồng quy có độ lớn lần lượt 3 $N$ và $4\,\text{N}$. Tính độ lớn hợp lực (kết quả theo đơn vị N).
\loigiai{
Xác định các đại lượng đã cho, chọn hệ thức phù hợp và thay số. Kết quả: $5\,\text{N}$.
}
\end{bt}

% ===== Câu 18 =====
% ID: AIQ_L10C3_FULL_0026
% Mức: VDC
\begin{bt}
% ID: AIQ_L10C3_FULL_0026
% Mức: VDC
Hai lực vuông góc đồng quy có độ lớn lần lượt 4 $N$ và $5\,\text{N}$. Tính độ lớn hợp lực (kết quả theo đơn vị N).
\loigiai{
Xác định các đại lượng đã cho, chọn hệ thức phù hợp và thay số. Kết quả: $6.4\,\text{N}$.
}
\end{bt}

% ===== Câu 19 =====
% ID: AIQ_L10C3_FULL_0027
% Mức: VDC
\begin{bt}
% ID: AIQ_L10C3_FULL_0027
% Mức: VDC
Hai lực vuông góc đồng quy có độ lớn lần lượt 5 $N$ và $6\,\text{N}$. Tính độ lớn hợp lực (kết quả theo đơn vị N).
\loigiai{
Xác định các đại lượng đã cho, chọn hệ thức phù hợp và thay số. Kết quả: $7.81\,\text{N}$.
}
\end{bt}

% ===== Câu 20 =====
% ID: AIQ_L10C3_FULL_0028
% Mức: NB
\dangbt{Tổng hợp hai lực đồng quy bằng quy tắc hình bình hành}
\begin{ex}
% ID: AIQ_L10C3_FULL_0028
% Mức: NB
Hai lực vuông góc đồng quy có độ lớn lần lượt 6 $N$ và $4\,\text{N}$. Tính độ lớn hợp lực (kết quả theo đơn vị N).
\choice
{\True 7.21}
{8.21}
{6.21}
{14.42}
\loigiai{
Áp dụng hệ thức phù hợp của bài học, thay số được $7.21\,\text{N}$.
}
\end{ex}

% ===== Câu 21 =====
% ID: AIQ_L10C3_FULL_0029
% Mức: NB
\begin{ex}
% ID: AIQ_L10C3_FULL_0029
% Mức: NB
Hai lực vuông góc đồng quy có độ lớn lần lượt 3 $N$ và $5\,\text{N}$. Tính độ lớn hợp lực (kết quả theo đơn vị N).
\choice
{6.83}
{\True 5.83}
{4.83}
{11.66}
\loigiai{
Áp dụng hệ thức phù hợp của bài học, thay số được $5.83\,\text{N}$.
}
\end{ex}

% ===== Câu 22 =====
% ID: AIQ_L10C3_FULL_0030
% Mức: NB
\begin{ex}
% ID: AIQ_L10C3_FULL_0030
% Mức: NB
Hai lực vuông góc đồng quy có độ lớn lần lượt 4 $N$ và $6\,\text{N}$. Tính độ lớn hợp lực (kết quả theo đơn vị N).
\choice
{6.21}
{8.21}
{\True 7.21}
{14.42}
\loigiai{
Áp dụng hệ thức phù hợp của bài học, thay số được $7.21\,\text{N}$.
}
\end{ex}

% ===== Câu 23 =====
% ID: AIQ_L10C3_FULL_0031
% Mức: NB
\begin{ex}
% ID: AIQ_L10C3_FULL_0031
% Mức: NB
Hai lực vuông góc đồng quy có độ lớn lần lượt 5 $N$ và $4\,\text{N}$. Tính độ lớn hợp lực (kết quả theo đơn vị N).
\choice
{12.8}
{7.4}
{5.4}
{\True 6.4}
\loigiai{
Áp dụng hệ thức phù hợp của bài học, thay số được $6.4\,\text{N}$.
}
\end{ex}

% ===== Câu 24 =====
% ID: AIQ_L10C3_FULL_0032
% Mức: TH
\begin{ex}
% ID: AIQ_L10C3_FULL_0032
% Mức: TH
Hai lực vuông góc đồng quy có độ lớn lần lượt 6 $N$ và $5\,\text{N}$. Tính độ lớn hợp lực (kết quả theo đơn vị N).
\choice
{\True 7.81}
{8.81}
{6.81}
{15.62}
\loigiai{
Áp dụng hệ thức phù hợp của bài học, thay số được $7.81\,\text{N}$.
}
\end{ex}

% ===== Câu 25 =====
% ID: AIQ_L10C3_FULL_0033
% Mức: TH
\begin{ex}
% ID: AIQ_L10C3_FULL_0033
% Mức: TH
Hai lực vuông góc đồng quy có độ lớn lần lượt 3 $N$ và $6\,\text{N}$. Tính độ lớn hợp lực (kết quả theo đơn vị N).
\choice
{7.71}
{\True 6.71}
{5.71}
{13.42}
\loigiai{
Áp dụng hệ thức phù hợp của bài học, thay số được $6.71\,\text{N}$.
}
\end{ex}

% ===== Câu 26 =====
% ID: AIQ_L10C3_FULL_0034
% Mức: TH
\begin{ex}
% ID: AIQ_L10C3_FULL_0034
% Mức: TH
Hai lực vuông góc đồng quy có độ lớn lần lượt 4 $N$ và $4\,\text{N}$. Tính độ lớn hợp lực (kết quả theo đơn vị N).
\choice
{4.66}
{6.66}
{\True 5.66}
{11.32}
\loigiai{
Áp dụng hệ thức phù hợp của bài học, thay số được $5.66\,\text{N}$.
}
\end{ex}

% ===== Câu 27 =====
% ID: AIQ_L10C3_FULL_0035
% Mức: TH
\begin{ex}
% ID: AIQ_L10C3_FULL_0035
% Mức: TH
Hai lực vuông góc đồng quy có độ lớn lần lượt 5 $N$ và $5\,\text{N}$. Tính độ lớn hợp lực (kết quả theo đơn vị N).
\choice
{14.14}
{8.07}
{6.07}
{\True 7.07}
\loigiai{
Áp dụng hệ thức phù hợp của bài học, thay số được $7.07\,\text{N}$.
}
\end{ex}

% ===== Câu 28 =====
% ID: AIQ_L10C3_FULL_0036
% Mức: VD
\begin{ex}
% ID: AIQ_L10C3_FULL_0036
% Mức: VD
Hai lực vuông góc đồng quy có độ lớn lần lượt 6 $N$ và $6\,\text{N}$. Tính độ lớn hợp lực (kết quả theo đơn vị N).
\choice
{\True 8.49}
{9.49}
{7.49}
{16.98}
\loigiai{
Áp dụng hệ thức phù hợp của bài học, thay số được $8.49\,\text{N}$.
}
\end{ex}

% ===== Câu 29 =====
% ID: AIQ_L10C3_FULL_0037
% Mức: VD
\begin{ex}
% ID: AIQ_L10C3_FULL_0037
% Mức: VD
Hai lực vuông góc đồng quy có độ lớn lần lượt 3 $N$ và $4\,\text{N}$. Tính độ lớn hợp lực (kết quả theo đơn vị N).
\choice
{6}
{\True 5}
{4}
{10}
\loigiai{
Áp dụng hệ thức phù hợp của bài học, thay số được $5\,\text{N}$.
}
\end{ex}

% ===== Câu 30 =====
% ID: AIQ_L10C3_FULL_0038
% Mức: TH
\begin{ex}
% ID: AIQ_L10C3_FULL_0038
% Mức: TH
Hai lực vuông góc đồng quy có độ lớn lần lượt 4 $N$ và $5\,\text{N}$. Tính độ lớn hợp lực (kết quả theo đơn vị N).
\choiceTF
{\True Giá trị cần tìm là $6.4\,\text{N}$.}
{Giá trị cần tìm là $7.4\,\text{N}$.}
{\True Phải xác định đầy đủ các lực hoặc đại lượng liên quan trước khi tính.}
{Có thể bỏ qua điều kiện áp dụng của công thức.}
\loigiai{
\begin{itemchoice}
\item (Đúng) Giá trị cần tìm là $6.4\,\text{N}$. (Vì): kết quả $6.4\,\text{N}$. b) Sai. c) Đúng. d) Sai vì phải kiểm tra điều kiện áp dụng
\item (Sai) Giá trị cần tìm là $7.4\,\text{N}$.
\item (Đúng) Phải xác định đầy đủ các lực hoặc đại lượng liên quan trước khi tính.
\item (Sai) Có thể bỏ qua điều kiện áp dụng của công thức.
\end{itemchoice}
}
\end{ex}

% ===== Câu 31 =====
% ID: AIQ_L10C3_FULL_0039
% Mức: TH
\begin{ex}
% ID: AIQ_L10C3_FULL_0039
% Mức: TH
Hai lực vuông góc đồng quy có độ lớn lần lượt 5 $N$ và $6\,\text{N}$. Tính độ lớn hợp lực (kết quả theo đơn vị N).
\choiceTF
{\True Giá trị cần tìm là $7.81\,\text{N}$.}
{Giá trị cần tìm là $8.81\,\text{N}$.}
{\True Phải xác định đầy đủ các lực hoặc đại lượng liên quan trước khi tính.}
{Có thể bỏ qua điều kiện áp dụng của công thức.}
\loigiai{
\begin{itemchoice}
\item (Đúng) Giá trị cần tìm là $7.81\,\text{N}$. (Vì): kết quả $7.81\,\text{N}$. b) Sai. c) Đúng. d) Sai vì phải kiểm tra điều kiện áp dụng
\item (Sai) Giá trị cần tìm là $8.81\,\text{N}$.
\item (Đúng) Phải xác định đầy đủ các lực hoặc đại lượng liên quan trước khi tính.
\item (Sai) Có thể bỏ qua điều kiện áp dụng của công thức.
\end{itemchoice}
}
\end{ex}

% ===== Câu 32 =====
% ID: AIQ_L10C3_FULL_0040
% Mức: VD
\begin{ex}
% ID: AIQ_L10C3_FULL_0040
% Mức: VD
Hai lực vuông góc đồng quy có độ lớn lần lượt 6 $N$ và $4\,\text{N}$. Tính độ lớn hợp lực (kết quả theo đơn vị N).
\choiceTF
{\True Giá trị cần tìm là $7.21\,\text{N}$.}
{Giá trị cần tìm là $8.21\,\text{N}$.}
{\True Phải xác định đầy đủ các lực hoặc đại lượng liên quan trước khi tính.}
{Có thể bỏ qua điều kiện áp dụng của công thức.}
\loigiai{
\begin{itemchoice}
\item (Đúng) Giá trị cần tìm là $7.21\,\text{N}$. (Vì): kết quả $7.21\,\text{N}$. b) Sai. c) Đúng. d) Sai vì phải kiểm tra điều kiện áp dụng
\item (Sai) Giá trị cần tìm là $8.21\,\text{N}$.
\item (Đúng) Phải xác định đầy đủ các lực hoặc đại lượng liên quan trước khi tính.
\item (Sai) Có thể bỏ qua điều kiện áp dụng của công thức.
\end{itemchoice}
}
\end{ex}

% ===== Câu 33 =====
% ID: AIQ_L10C3_FULL_0041
% Mức: VD
\begin{ex}
% ID: AIQ_L10C3_FULL_0041
% Mức: VD
Hai lực vuông góc đồng quy có độ lớn lần lượt 3 $N$ và $5\,\text{N}$. Tính độ lớn hợp lực (kết quả theo đơn vị N).
\choiceTF
{\True Giá trị cần tìm là $5.83\,\text{N}$.}
{Giá trị cần tìm là $6.83\,\text{N}$.}
{\True Phải xác định đầy đủ các lực hoặc đại lượng liên quan trước khi tính.}
{Có thể bỏ qua điều kiện áp dụng của công thức.}
\loigiai{
\begin{itemchoice}
\item (Đúng) Giá trị cần tìm là $5.83\,\text{N}$. (Vì): kết quả $5.83\,\text{N}$. b) Sai. c) Đúng. d) Sai vì phải kiểm tra điều kiện áp dụng
\item (Sai) Giá trị cần tìm là $6.83\,\text{N}$.
\item (Đúng) Phải xác định đầy đủ các lực hoặc đại lượng liên quan trước khi tính.
\item (Sai) Có thể bỏ qua điều kiện áp dụng của công thức.
\end{itemchoice}
}
\end{ex}

% ===== Câu 34 =====
% ID: AIQ_L10C3_FULL_0042
% Mức: VD
\begin{ex}
% ID: AIQ_L10C3_FULL_0042
% Mức: VD
Hai lực vuông góc đồng quy có độ lớn lần lượt 4 $N$ và $6\,\text{N}$. Tính độ lớn hợp lực (kết quả theo đơn vị N).
\choiceTF
{\True Giá trị cần tìm là $7.21\,\text{N}$.}
{Giá trị cần tìm là $8.21\,\text{N}$.}
{\True Phải xác định đầy đủ các lực hoặc đại lượng liên quan trước khi tính.}
{Có thể bỏ qua điều kiện áp dụng của công thức.}
\loigiai{
\begin{itemchoice}
\item (Đúng) Giá trị cần tìm là $7.21\,\text{N}$. (Vì): kết quả $7.21\,\text{N}$. b) Sai. c) Đúng. d) Sai vì phải kiểm tra điều kiện áp dụng
\item (Sai) Giá trị cần tìm là $8.21\,\text{N}$.
\item (Đúng) Phải xác định đầy đủ các lực hoặc đại lượng liên quan trước khi tính.
\item (Sai) Có thể bỏ qua điều kiện áp dụng của công thức.
\end{itemchoice}
}
\end{ex}

% ===== Câu 35 =====
% ID: AIQ_L10C3_FULL_0043
% Mức: TH
\begin{ex}
% ID: AIQ_L10C3_FULL_0043
% Mức: TH
Hai lực vuông góc đồng quy có độ lớn lần lượt 5 $N$ và $4\,\text{N}$. Tính độ lớn hợp lực (kết quả theo đơn vị N).
\shortans{6.4}
\loigiai{
Thay số vào hệ thức của bài học thu được $6.4\,\text{N}$.
}
\end{ex}

% ===== Câu 36 =====
% ID: AIQ_L10C3_FULL_0044
% Mức: TH
\begin{ex}
% ID: AIQ_L10C3_FULL_0044
% Mức: TH
Hai lực vuông góc đồng quy có độ lớn lần lượt 6 $N$ và $5\,\text{N}$. Tính độ lớn hợp lực (kết quả theo đơn vị N).
\shortans{7.81}
\loigiai{
Thay số vào hệ thức của bài học thu được $7.81\,\text{N}$.
}
\end{ex}

% ===== Câu 37 =====
% ID: AIQ_L10C3_FULL_0045
% Mức: TH
\begin{ex}
% ID: AIQ_L10C3_FULL_0045
% Mức: TH
Hai lực vuông góc đồng quy có độ lớn lần lượt 3 $N$ và $6\,\text{N}$. Tính độ lớn hợp lực (kết quả theo đơn vị N).
\shortans{6.71}
\loigiai{
Thay số vào hệ thức của bài học thu được $6.71\,\text{N}$.
}
\end{ex}

% ===== Câu 38 =====
% ID: AIQ_L10C3_FULL_0046
% Mức: VD
\begin{ex}
% ID: AIQ_L10C3_FULL_0046
% Mức: VD
Hai lực vuông góc đồng quy có độ lớn lần lượt 4 $N$ và $4\,\text{N}$. Tính độ lớn hợp lực (kết quả theo đơn vị N).
\shortans{5.66}
\loigiai{
Thay số vào hệ thức của bài học thu được $5.66\,\text{N}$.
}
\end{ex}

% ===== Câu 39 =====
% ID: AIQ_L10C3_FULL_0047
% Mức: VD
\begin{ex}
% ID: AIQ_L10C3_FULL_0047
% Mức: VD
Hai lực vuông góc đồng quy có độ lớn lần lượt 5 $N$ và $5\,\text{N}$. Tính độ lớn hợp lực (kết quả theo đơn vị N).
\shortans{7.07}
\loigiai{
Thay số vào hệ thức của bài học thu được $7.07\,\text{N}$.
}
\end{ex}

% ===== Câu 40 =====
% ID: AIQ_L10C3_FULL_0048
% Mức: VD
\begin{ex}
% ID: AIQ_L10C3_FULL_0048
% Mức: VD
Hai lực vuông góc đồng quy có độ lớn lần lượt 6 $N$ và $6\,\text{N}$. Tính độ lớn hợp lực (kết quả theo đơn vị N).
\shortans{8.49}
\loigiai{
Thay số vào hệ thức của bài học thu được $8.49\,\text{N}$.
}
\end{ex}

% ===== Câu 41 =====
% ID: AIQ_L10C3_FULL_0049
% Mức: VD
\begin{bt}
% ID: AIQ_L10C3_FULL_0049
% Mức: VD
Hai lực vuông góc đồng quy có độ lớn lần lượt 3 $N$ và $4\,\text{N}$. Tính độ lớn hợp lực (kết quả theo đơn vị N).
\loigiai{
Xác định các đại lượng đã cho, chọn hệ thức phù hợp và thay số. Kết quả: $5\,\text{N}$.
}
\end{bt}

% ===== Câu 42 =====
% ID: AIQ_L10C3_FULL_0050
% Mức: VD
\begin{bt}
% ID: AIQ_L10C3_FULL_0050
% Mức: VD
Hai lực vuông góc đồng quy có độ lớn lần lượt 4 $N$ và $5\,\text{N}$. Tính độ lớn hợp lực (kết quả theo đơn vị N).
\loigiai{
Xác định các đại lượng đã cho, chọn hệ thức phù hợp và thay số. Kết quả: $6.4\,\text{N}$.
}
\end{bt}

% ===== Câu 43 =====
% ID: AIQ_L10C3_FULL_0051
% Mức: VD
\begin{bt}
% ID: AIQ_L10C3_FULL_0051
% Mức: VD
Hai lực vuông góc đồng quy có độ lớn lần lượt 5 $N$ và $6\,\text{N}$. Tính độ lớn hợp lực (kết quả theo đơn vị N).
\loigiai{
Xác định các đại lượng đã cho, chọn hệ thức phù hợp và thay số. Kết quả: $7.81\,\text{N}$.
}
\end{bt}

% ===== Câu 44 =====
% ID: AIQ_L10C3_FULL_0052
% Mức: VDC
\begin{bt}
% ID: AIQ_L10C3_FULL_0052
% Mức: VDC
Hai lực vuông góc đồng quy có độ lớn lần lượt 6 $N$ và $4\,\text{N}$. Tính độ lớn hợp lực (kết quả theo đơn vị N).
\loigiai{
Xác định các đại lượng đã cho, chọn hệ thức phù hợp và thay số. Kết quả: $7.21\,\text{N}$.
}
\end{bt}

% ===== Câu 45 =====
% ID: AIQ_L10C3_FULL_0053
% Mức: VDC
\begin{bt}
% ID: AIQ_L10C3_FULL_0053
% Mức: VDC
Hai lực vuông góc đồng quy có độ lớn lần lượt 3 $N$ và $5\,\text{N}$. Tính độ lớn hợp lực (kết quả theo đơn vị N).
\loigiai{
Xác định các đại lượng đã cho, chọn hệ thức phù hợp và thay số. Kết quả: $5.83\,\text{N}$.
}
\end{bt}

% ===== Câu 46 =====
% ID: AIQ_L10C3_FULL_0054
% Mức: VDC
\begin{bt}
% ID: AIQ_L10C3_FULL_0054
% Mức: VDC
Hai lực vuông góc đồng quy có độ lớn lần lượt 4 $N$ và $6\,\text{N}$. Tính độ lớn hợp lực (kết quả theo đơn vị N).
\loigiai{
Xác định các đại lượng đã cho, chọn hệ thức phù hợp và thay số. Kết quả: $7.21\,\text{N}$.
}
\end{bt}

% ===== Câu 47 =====
% ID: AIQ_L10C3_FULL_0055
% Mức: NB
\dangbt{Phân tích một lực thành hai lực thành phần}
\begin{ex}
% ID: AIQ_L10C3_FULL_0055
% Mức: NB
Hai lực vuông góc đồng quy có độ lớn lần lượt 5 $N$ và $4\,\text{N}$. Tính độ lớn hợp lực (kết quả theo đơn vị N).
\choice
{\True 6.4}
{7.4}
{5.4}
{12.8}
\loigiai{
Áp dụng hệ thức phù hợp của bài học, thay số được $6.4\,\text{N}$.
}
\end{ex}

% ===== Câu 48 =====
% ID: AIQ_L10C3_FULL_0056
% Mức: NB
\begin{ex}
% ID: AIQ_L10C3_FULL_0056
% Mức: NB
Hai lực vuông góc đồng quy có độ lớn lần lượt 6 $N$ và $5\,\text{N}$. Tính độ lớn hợp lực (kết quả theo đơn vị N).
\choice
{8.81}
{\True 7.81}
{6.81}
{15.62}
\loigiai{
Áp dụng hệ thức phù hợp của bài học, thay số được $7.81\,\text{N}$.
}
\end{ex}

% ===== Câu 49 =====
% ID: AIQ_L10C3_FULL_0057
% Mức: NB
\begin{ex}
% ID: AIQ_L10C3_FULL_0057
% Mức: NB
Hai lực vuông góc đồng quy có độ lớn lần lượt 3 $N$ và $6\,\text{N}$. Tính độ lớn hợp lực (kết quả theo đơn vị N).
\choice
{5.71}
{7.71}
{\True 6.71}
{13.42}
\loigiai{
Áp dụng hệ thức phù hợp của bài học, thay số được $6.71\,\text{N}$.
}
\end{ex}

% ===== Câu 50 =====
% ID: AIQ_L10C3_FULL_0058
% Mức: NB
\begin{ex}
% ID: AIQ_L10C3_FULL_0058
% Mức: NB
Hai lực vuông góc đồng quy có độ lớn lần lượt 4 $N$ và $4\,\text{N}$. Tính độ lớn hợp lực (kết quả theo đơn vị N).
\choice
{11.32}
{6.66}
{4.66}
{\True 5.66}
\loigiai{
Áp dụng hệ thức phù hợp của bài học, thay số được $5.66\,\text{N}$.
}
\end{ex}

% ===== Câu 51 =====
% ID: AIQ_L10C3_FULL_0059
% Mức: TH
\begin{ex}
% ID: AIQ_L10C3_FULL_0059
% Mức: TH
Hai lực vuông góc đồng quy có độ lớn lần lượt 5 $N$ và $5\,\text{N}$. Tính độ lớn hợp lực (kết quả theo đơn vị N).
\choice
{\True 7.07}
{8.07}
{6.07}
{14.14}
\loigiai{
Áp dụng hệ thức phù hợp của bài học, thay số được $7.07\,\text{N}$.
}
\end{ex}

% ===== Câu 52 =====
% ID: AIQ_L10C3_FULL_0060
% Mức: TH
\begin{ex}
% ID: AIQ_L10C3_FULL_0060
% Mức: TH
Hai lực vuông góc đồng quy có độ lớn lần lượt 6 $N$ và $6\,\text{N}$. Tính độ lớn hợp lực (kết quả theo đơn vị N).
\choice
{9.49}
{\True 8.49}
{7.49}
{16.98}
\loigiai{
Áp dụng hệ thức phù hợp của bài học, thay số được $8.49\,\text{N}$.
}
\end{ex}

% ===== Câu 53 =====
% ID: AIQ_L10C3_FULL_0061
% Mức: TH
\begin{ex}
% ID: AIQ_L10C3_FULL_0061
% Mức: TH
Hai lực vuông góc đồng quy có độ lớn lần lượt 3 $N$ và $4\,\text{N}$. Tính độ lớn hợp lực (kết quả theo đơn vị N).
\choice
{4}
{6}
{\True 5}
{10}
\loigiai{
Áp dụng hệ thức phù hợp của bài học, thay số được $5\,\text{N}$.
}
\end{ex}

% ===== Câu 54 =====
% ID: AIQ_L10C3_FULL_0062
% Mức: TH
\begin{ex}
% ID: AIQ_L10C3_FULL_0062
% Mức: TH
Hai lực vuông góc đồng quy có độ lớn lần lượt 4 $N$ và $5\,\text{N}$. Tính độ lớn hợp lực (kết quả theo đơn vị N).
\choice
{12.8}
{7.4}
{5.4}
{\True 6.4}
\loigiai{
Áp dụng hệ thức phù hợp của bài học, thay số được $6.4\,\text{N}$.
}
\end{ex}

% ===== Câu 55 =====
% ID: AIQ_L10C3_FULL_0063
% Mức: VD
\begin{ex}
% ID: AIQ_L10C3_FULL_0063
% Mức: VD
Hai lực vuông góc đồng quy có độ lớn lần lượt 5 $N$ và $6\,\text{N}$. Tính độ lớn hợp lực (kết quả theo đơn vị N).
\choice
{\True 7.81}
{8.81}
{6.81}
{15.62}
\loigiai{
Áp dụng hệ thức phù hợp của bài học, thay số được $7.81\,\text{N}$.
}
\end{ex}

% ===== Câu 56 =====
% ID: AIQ_L10C3_FULL_0064
% Mức: VD
\begin{ex}
% ID: AIQ_L10C3_FULL_0064
% Mức: VD
Hai lực vuông góc đồng quy có độ lớn lần lượt 6 $N$ và $4\,\text{N}$. Tính độ lớn hợp lực (kết quả theo đơn vị N).
\choice
{8.21}
{\True 7.21}
{6.21}
{14.42}
\loigiai{
Áp dụng hệ thức phù hợp của bài học, thay số được $7.21\,\text{N}$.
}
\end{ex}

% ===== Câu 57 =====
% ID: AIQ_L10C3_FULL_0065
% Mức: TH
\begin{ex}
% ID: AIQ_L10C3_FULL_0065
% Mức: TH
Hai lực vuông góc đồng quy có độ lớn lần lượt 3 $N$ và $5\,\text{N}$. Tính độ lớn hợp lực (kết quả theo đơn vị N).
\choiceTF
{\True Giá trị cần tìm là $5.83\,\text{N}$.}
{Giá trị cần tìm là $6.83\,\text{N}$.}
{\True Phải xác định đầy đủ các lực hoặc đại lượng liên quan trước khi tính.}
{Có thể bỏ qua điều kiện áp dụng của công thức.}
\loigiai{
\begin{itemchoice}
\item (Đúng) Giá trị cần tìm là $5.83\,\text{N}$. (Vì): kết quả $5.83\,\text{N}$. b) Sai. c) Đúng. d) Sai vì phải kiểm tra điều kiện áp dụng
\item (Sai) Giá trị cần tìm là $6.83\,\text{N}$.
\item (Đúng) Phải xác định đầy đủ các lực hoặc đại lượng liên quan trước khi tính.
\item (Sai) Có thể bỏ qua điều kiện áp dụng của công thức.
\end{itemchoice}
}
\end{ex}

% ===== Câu 58 =====
% ID: AIQ_L10C3_FULL_0066
% Mức: TH
\begin{ex}
% ID: AIQ_L10C3_FULL_0066
% Mức: TH
Hai lực vuông góc đồng quy có độ lớn lần lượt 4 $N$ và $6\,\text{N}$. Tính độ lớn hợp lực (kết quả theo đơn vị N).
\choiceTF
{\True Giá trị cần tìm là $7.21\,\text{N}$.}
{Giá trị cần tìm là $8.21\,\text{N}$.}
{\True Phải xác định đầy đủ các lực hoặc đại lượng liên quan trước khi tính.}
{Có thể bỏ qua điều kiện áp dụng của công thức.}
\loigiai{
\begin{itemchoice}
\item (Đúng) Giá trị cần tìm là $7.21\,\text{N}$. (Vì): kết quả $7.21\,\text{N}$. b) Sai. c) Đúng. d) Sai vì phải kiểm tra điều kiện áp dụng
\item (Sai) Giá trị cần tìm là $8.21\,\text{N}$.
\item (Đúng) Phải xác định đầy đủ các lực hoặc đại lượng liên quan trước khi tính.
\item (Sai) Có thể bỏ qua điều kiện áp dụng của công thức.
\end{itemchoice}
}
\end{ex}

% ===== Câu 59 =====
% ID: AIQ_L10C3_FULL_0067
% Mức: VD
\begin{ex}
% ID: AIQ_L10C3_FULL_0067
% Mức: VD
Hai lực vuông góc đồng quy có độ lớn lần lượt 5 $N$ và $4\,\text{N}$. Tính độ lớn hợp lực (kết quả theo đơn vị N).
\choiceTF
{\True Giá trị cần tìm là $6.4\,\text{N}$.}
{Giá trị cần tìm là $7.4\,\text{N}$.}
{\True Phải xác định đầy đủ các lực hoặc đại lượng liên quan trước khi tính.}
{Có thể bỏ qua điều kiện áp dụng của công thức.}
\loigiai{
\begin{itemchoice}
\item (Đúng) Giá trị cần tìm là $6.4\,\text{N}$. (Vì): kết quả $6.4\,\text{N}$. b) Sai. c) Đúng. d) Sai vì phải kiểm tra điều kiện áp dụng
\item (Sai) Giá trị cần tìm là $7.4\,\text{N}$.
\item (Đúng) Phải xác định đầy đủ các lực hoặc đại lượng liên quan trước khi tính.
\item (Sai) Có thể bỏ qua điều kiện áp dụng của công thức.
\end{itemchoice}
}
\end{ex}

% ===== Câu 60 =====
% ID: AIQ_L10C3_FULL_0068
% Mức: VD
\begin{ex}
% ID: AIQ_L10C3_FULL_0068
% Mức: VD
Hai lực vuông góc đồng quy có độ lớn lần lượt 6 $N$ và $5\,\text{N}$. Tính độ lớn hợp lực (kết quả theo đơn vị N).
\choiceTF
{\True Giá trị cần tìm là $7.81\,\text{N}$.}
{Giá trị cần tìm là $8.81\,\text{N}$.}
{\True Phải xác định đầy đủ các lực hoặc đại lượng liên quan trước khi tính.}
{Có thể bỏ qua điều kiện áp dụng của công thức.}
\loigiai{
\begin{itemchoice}
\item (Đúng) Giá trị cần tìm là $7.81\,\text{N}$. (Vì): kết quả $7.81\,\text{N}$. b) Sai. c) Đúng. d) Sai vì phải kiểm tra điều kiện áp dụng
\item (Sai) Giá trị cần tìm là $8.81\,\text{N}$.
\item (Đúng) Phải xác định đầy đủ các lực hoặc đại lượng liên quan trước khi tính.
\item (Sai) Có thể bỏ qua điều kiện áp dụng của công thức.
\end{itemchoice}
}
\end{ex}

% ===== Câu 61 =====
% ID: AIQ_L10C3_FULL_0069
% Mức: VD
\begin{ex}
% ID: AIQ_L10C3_FULL_0069
% Mức: VD
Hai lực vuông góc đồng quy có độ lớn lần lượt 3 $N$ và $6\,\text{N}$. Tính độ lớn hợp lực (kết quả theo đơn vị N).
\choiceTF
{\True Giá trị cần tìm là $6.71\,\text{N}$.}
{Giá trị cần tìm là $7.71\,\text{N}$.}
{\True Phải xác định đầy đủ các lực hoặc đại lượng liên quan trước khi tính.}
{Có thể bỏ qua điều kiện áp dụng của công thức.}
\loigiai{
\begin{itemchoice}
\item (Đúng) Giá trị cần tìm là $6.71\,\text{N}$. (Vì): kết quả $6.71\,\text{N}$. b) Sai. c) Đúng. d) Sai vì phải kiểm tra điều kiện áp dụng
\item (Sai) Giá trị cần tìm là $7.71\,\text{N}$.
\item (Đúng) Phải xác định đầy đủ các lực hoặc đại lượng liên quan trước khi tính.
\item (Sai) Có thể bỏ qua điều kiện áp dụng của công thức.
\end{itemchoice}
}
\end{ex}

% ===== Câu 62 =====
% ID: AIQ_L10C3_FULL_0070
% Mức: TH
\begin{ex}
% ID: AIQ_L10C3_FULL_0070
% Mức: TH
Hai lực vuông góc đồng quy có độ lớn lần lượt 4 $N$ và $4\,\text{N}$. Tính độ lớn hợp lực (kết quả theo đơn vị N).
\shortans{5.66}
\loigiai{
Thay số vào hệ thức của bài học thu được $5.66\,\text{N}$.
}
\end{ex}

% ===== Câu 63 =====
% ID: AIQ_L10C3_FULL_0071
% Mức: TH
\begin{ex}
% ID: AIQ_L10C3_FULL_0071
% Mức: TH
Hai lực vuông góc đồng quy có độ lớn lần lượt 5 $N$ và $5\,\text{N}$. Tính độ lớn hợp lực (kết quả theo đơn vị N).
\shortans{7.07}
\loigiai{
Thay số vào hệ thức của bài học thu được $7.07\,\text{N}$.
}
\end{ex}

% ===== Câu 64 =====
% ID: AIQ_L10C3_FULL_0072
% Mức: TH
\begin{ex}
% ID: AIQ_L10C3_FULL_0072
% Mức: TH
Hai lực vuông góc đồng quy có độ lớn lần lượt 6 $N$ và $6\,\text{N}$. Tính độ lớn hợp lực (kết quả theo đơn vị N).
\shortans{8.49}
\loigiai{
Thay số vào hệ thức của bài học thu được $8.49\,\text{N}$.
}
\end{ex}

% ===== Câu 65 =====
% ID: AIQ_L10C3_FULL_0073
% Mức: VD
\begin{ex}
% ID: AIQ_L10C3_FULL_0073
% Mức: VD
Hai lực vuông góc đồng quy có độ lớn lần lượt 3 $N$ và $4\,\text{N}$. Tính độ lớn hợp lực (kết quả theo đơn vị N).
\shortans{5}
\loigiai{
Thay số vào hệ thức của bài học thu được $5\,\text{N}$.
}
\end{ex}

% ===== Câu 66 =====
% ID: AIQ_L10C3_FULL_0074
% Mức: VD
\begin{ex}
% ID: AIQ_L10C3_FULL_0074
% Mức: VD
Hai lực vuông góc đồng quy có độ lớn lần lượt 4 $N$ và $5\,\text{N}$. Tính độ lớn hợp lực (kết quả theo đơn vị N).
\shortans{6.4}
\loigiai{
Thay số vào hệ thức của bài học thu được $6.4\,\text{N}$.
}
\end{ex}

% ===== Câu 67 =====
% ID: AIQ_L10C3_FULL_0075
% Mức: VD
\begin{ex}
% ID: AIQ_L10C3_FULL_0075
% Mức: VD
Hai lực vuông góc đồng quy có độ lớn lần lượt 5 $N$ và $6\,\text{N}$. Tính độ lớn hợp lực (kết quả theo đơn vị N).
\shortans{7.81}
\loigiai{
Thay số vào hệ thức của bài học thu được $7.81\,\text{N}$.
}
\end{ex}

% ===== Câu 68 =====
% ID: AIQ_L10C3_FULL_0076
% Mức: VD
\begin{bt}
% ID: AIQ_L10C3_FULL_0076
% Mức: VD
Hai lực vuông góc đồng quy có độ lớn lần lượt 6 $N$ và $4\,\text{N}$. Tính độ lớn hợp lực (kết quả theo đơn vị N).
\loigiai{
Xác định các đại lượng đã cho, chọn hệ thức phù hợp và thay số. Kết quả: $7.21\,\text{N}$.
}
\end{bt}

% ===== Câu 69 =====
% ID: AIQ_L10C3_FULL_0077
% Mức: VD
\begin{bt}
% ID: AIQ_L10C3_FULL_0077
% Mức: VD
Hai lực vuông góc đồng quy có độ lớn lần lượt 3 $N$ và $5\,\text{N}$. Tính độ lớn hợp lực (kết quả theo đơn vị N).
\loigiai{
Xác định các đại lượng đã cho, chọn hệ thức phù hợp và thay số. Kết quả: $5.83\,\text{N}$.
}
\end{bt}

% ===== Câu 70 =====
% ID: AIQ_L10C3_FULL_0078
% Mức: VD
\begin{bt}
% ID: AIQ_L10C3_FULL_0078
% Mức: VD
Hai lực vuông góc đồng quy có độ lớn lần lượt 4 $N$ và $6\,\text{N}$. Tính độ lớn hợp lực (kết quả theo đơn vị N).
\loigiai{
Xác định các đại lượng đã cho, chọn hệ thức phù hợp và thay số. Kết quả: $7.21\,\text{N}$.
}
\end{bt}

% ===== Câu 71 =====
% ID: AIQ_L10C3_FULL_0079
% Mức: VDC
\begin{bt}
% ID: AIQ_L10C3_FULL_0079
% Mức: VDC
Hai lực vuông góc đồng quy có độ lớn lần lượt 5 $N$ và $4\,\text{N}$. Tính độ lớn hợp lực (kết quả theo đơn vị N).
\loigiai{
Xác định các đại lượng đã cho, chọn hệ thức phù hợp và thay số. Kết quả: $6.4\,\text{N}$.
}
\end{bt}

% ===== Câu 72 =====
% ID: AIQ_L10C3_FULL_0080
% Mức: VDC
\begin{bt}
% ID: AIQ_L10C3_FULL_0080
% Mức: VDC
Hai lực vuông góc đồng quy có độ lớn lần lượt 6 $N$ và $5\,\text{N}$. Tính độ lớn hợp lực (kết quả theo đơn vị N).
\loigiai{
Xác định các đại lượng đã cho, chọn hệ thức phù hợp và thay số. Kết quả: $7.81\,\text{N}$.
}
\end{bt}

% ===== Câu 73 =====
% ID: AIQ_L10C3_FULL_0081
% Mức: VDC
\begin{bt}
% ID: AIQ_L10C3_FULL_0081
% Mức: VDC
Hai lực vuông góc đồng quy có độ lớn lần lượt 3 $N$ và $6\,\text{N}$. Tính độ lớn hợp lực (kết quả theo đơn vị N).
\loigiai{
Xác định các đại lượng đã cho, chọn hệ thức phù hợp và thay số. Kết quả: $6.71\,\text{N}$.
}
\end{bt}

% ===== Câu 74 =====
% ID: AIQ_L10C3_FULL_0082
% Mức: NB
\dangbt{Xác định hợp lực của nhiều lực đồng quy}
\begin{ex}
% ID: AIQ_L10C3_FULL_0082
% Mức: NB
Hai lực vuông góc đồng quy có độ lớn lần lượt 4 $N$ và $4\,\text{N}$. Tính độ lớn hợp lực (kết quả theo đơn vị N).
\choice
{\True 5.66}
{6.66}
{4.66}
{11.32}
\loigiai{
Áp dụng hệ thức phù hợp của bài học, thay số được $5.66\,\text{N}$.
}
\end{ex}

% ===== Câu 75 =====
% ID: AIQ_L10C3_FULL_0083
% Mức: NB
\begin{ex}
% ID: AIQ_L10C3_FULL_0083
% Mức: NB
Hai lực vuông góc đồng quy có độ lớn lần lượt 5 $N$ và $5\,\text{N}$. Tính độ lớn hợp lực (kết quả theo đơn vị N).
\choice
{8.07}
{\True 7.07}
{6.07}
{14.14}
\loigiai{
Áp dụng hệ thức phù hợp của bài học, thay số được $7.07\,\text{N}$.
}
\end{ex}

% ===== Câu 76 =====
% ID: AIQ_L10C3_FULL_0084
% Mức: NB
\begin{ex}
% ID: AIQ_L10C3_FULL_0084
% Mức: NB
Hai lực vuông góc đồng quy có độ lớn lần lượt 6 $N$ và $6\,\text{N}$. Tính độ lớn hợp lực (kết quả theo đơn vị N).
\choice
{7.49}
{9.49}
{\True 8.49}
{16.98}
\loigiai{
Áp dụng hệ thức phù hợp của bài học, thay số được $8.49\,\text{N}$.
}
\end{ex}

% ===== Câu 77 =====
% ID: AIQ_L10C3_FULL_0085
% Mức: NB
\begin{ex}
% ID: AIQ_L10C3_FULL_0085
% Mức: NB
Hai lực vuông góc đồng quy có độ lớn lần lượt 3 $N$ và $4\,\text{N}$. Tính độ lớn hợp lực (kết quả theo đơn vị N).
\choice
{10}
{6}
{4}
{\True 5}
\loigiai{
Áp dụng hệ thức phù hợp của bài học, thay số được $5\,\text{N}$.
}
\end{ex}

% ===== Câu 78 =====
% ID: AIQ_L10C3_FULL_0086
% Mức: TH
\begin{ex}
% ID: AIQ_L10C3_FULL_0086
% Mức: TH
Hai lực vuông góc đồng quy có độ lớn lần lượt 4 $N$ và $5\,\text{N}$. Tính độ lớn hợp lực (kết quả theo đơn vị N).
\choice
{\True 6.4}
{7.4}
{5.4}
{12.8}
\loigiai{
Áp dụng hệ thức phù hợp của bài học, thay số được $6.4\,\text{N}$.
}
\end{ex}

% ===== Câu 79 =====
% ID: AIQ_L10C3_FULL_0087
% Mức: TH
\begin{ex}
% ID: AIQ_L10C3_FULL_0087
% Mức: TH
Hai lực vuông góc đồng quy có độ lớn lần lượt 5 $N$ và $6\,\text{N}$. Tính độ lớn hợp lực (kết quả theo đơn vị N).
\choice
{8.81}
{\True 7.81}
{6.81}
{15.62}
\loigiai{
Áp dụng hệ thức phù hợp của bài học, thay số được $7.81\,\text{N}$.
}
\end{ex}

% ===== Câu 80 =====
% ID: AIQ_L10C3_FULL_0088
% Mức: TH
\begin{ex}
% ID: AIQ_L10C3_FULL_0088
% Mức: TH
Hai lực vuông góc đồng quy có độ lớn lần lượt 6 $N$ và $4\,\text{N}$. Tính độ lớn hợp lực (kết quả theo đơn vị N).
\choice
{6.21}
{8.21}
{\True 7.21}
{14.42}
\loigiai{
Áp dụng hệ thức phù hợp của bài học, thay số được $7.21\,\text{N}$.
}
\end{ex}

% ===== Câu 81 =====
% ID: AIQ_L10C3_FULL_0089
% Mức: TH
\begin{ex}
% ID: AIQ_L10C3_FULL_0089
% Mức: TH
Hai lực vuông góc đồng quy có độ lớn lần lượt 3 $N$ và $5\,\text{N}$. Tính độ lớn hợp lực (kết quả theo đơn vị N).
\choice
{11.66}
{6.83}
{4.83}
{\True 5.83}
\loigiai{
Áp dụng hệ thức phù hợp của bài học, thay số được $5.83\,\text{N}$.
}
\end{ex}

% ===== Câu 82 =====
% ID: AIQ_L10C3_FULL_0090
% Mức: VD
\begin{ex}
% ID: AIQ_L10C3_FULL_0090
% Mức: VD
Hai lực vuông góc đồng quy có độ lớn lần lượt 4 $N$ và $6\,\text{N}$. Tính độ lớn hợp lực (kết quả theo đơn vị N).
\choice
{\True 7.21}
{8.21}
{6.21}
{14.42}
\loigiai{
Áp dụng hệ thức phù hợp của bài học, thay số được $7.21\,\text{N}$.
}
\end{ex}

% ===== Câu 83 =====
% ID: AIQ_L10C3_FULL_0091
% Mức: VD
\begin{ex}
% ID: AIQ_L10C3_FULL_0091
% Mức: VD
Hai lực vuông góc đồng quy có độ lớn lần lượt 5 $N$ và $4\,\text{N}$. Tính độ lớn hợp lực (kết quả theo đơn vị N).
\choice
{7.4}
{\True 6.4}
{5.4}
{12.8}
\loigiai{
Áp dụng hệ thức phù hợp của bài học, thay số được $6.4\,\text{N}$.
}
\end{ex}

% ===== Câu 84 =====
% ID: AIQ_L10C3_FULL_0092
% Mức: TH
\begin{ex}
% ID: AIQ_L10C3_FULL_0092
% Mức: TH
Hai lực vuông góc đồng quy có độ lớn lần lượt 6 $N$ và $5\,\text{N}$. Tính độ lớn hợp lực (kết quả theo đơn vị N).
\choiceTF
{\True Giá trị cần tìm là $7.81\,\text{N}$.}
{Giá trị cần tìm là $8.81\,\text{N}$.}
{\True Phải xác định đầy đủ các lực hoặc đại lượng liên quan trước khi tính.}
{Có thể bỏ qua điều kiện áp dụng của công thức.}
\loigiai{
\begin{itemchoice}
\item (Đúng) Giá trị cần tìm là $7.81\,\text{N}$. (Vì): kết quả $7.81\,\text{N}$. b) Sai. c) Đúng. d) Sai vì phải kiểm tra điều kiện áp dụng
\item (Sai) Giá trị cần tìm là $8.81\,\text{N}$.
\item (Đúng) Phải xác định đầy đủ các lực hoặc đại lượng liên quan trước khi tính.
\item (Sai) Có thể bỏ qua điều kiện áp dụng của công thức.
\end{itemchoice}
}
\end{ex}

% ===== Câu 85 =====
% ID: AIQ_L10C3_FULL_0093
% Mức: TH
\begin{ex}
% ID: AIQ_L10C3_FULL_0093
% Mức: TH
Hai lực vuông góc đồng quy có độ lớn lần lượt 3 $N$ và $6\,\text{N}$. Tính độ lớn hợp lực (kết quả theo đơn vị N).
\choiceTF
{\True Giá trị cần tìm là $6.71\,\text{N}$.}
{Giá trị cần tìm là $7.71\,\text{N}$.}
{\True Phải xác định đầy đủ các lực hoặc đại lượng liên quan trước khi tính.}
{Có thể bỏ qua điều kiện áp dụng của công thức.}
\loigiai{
\begin{itemchoice}
\item (Đúng) Giá trị cần tìm là $6.71\,\text{N}$. (Vì): kết quả $6.71\,\text{N}$. b) Sai. c) Đúng. d) Sai vì phải kiểm tra điều kiện áp dụng
\item (Sai) Giá trị cần tìm là $7.71\,\text{N}$.
\item (Đúng) Phải xác định đầy đủ các lực hoặc đại lượng liên quan trước khi tính.
\item (Sai) Có thể bỏ qua điều kiện áp dụng của công thức.
\end{itemchoice}
}
\end{ex}

% ===== Câu 86 =====
% ID: AIQ_L10C3_FULL_0094
% Mức: VD
\begin{ex}
% ID: AIQ_L10C3_FULL_0094
% Mức: VD
Hai lực vuông góc đồng quy có độ lớn lần lượt 4 $N$ và $4\,\text{N}$. Tính độ lớn hợp lực (kết quả theo đơn vị N).
\choiceTF
{\True Giá trị cần tìm là $5.66\,\text{N}$.}
{Giá trị cần tìm là $6.66\,\text{N}$.}
{\True Phải xác định đầy đủ các lực hoặc đại lượng liên quan trước khi tính.}
{Có thể bỏ qua điều kiện áp dụng của công thức.}
\loigiai{
\begin{itemchoice}
\item (Đúng) Giá trị cần tìm là $5.66\,\text{N}$. (Vì): kết quả $5.66\,\text{N}$. b) Sai. c) Đúng. d) Sai vì phải kiểm tra điều kiện áp dụng
\item (Sai) Giá trị cần tìm là $6.66\,\text{N}$.
\item (Đúng) Phải xác định đầy đủ các lực hoặc đại lượng liên quan trước khi tính.
\item (Sai) Có thể bỏ qua điều kiện áp dụng của công thức.
\end{itemchoice}
}
\end{ex}

% ===== Câu 87 =====
% ID: AIQ_L10C3_FULL_0095
% Mức: VD
\begin{ex}
% ID: AIQ_L10C3_FULL_0095
% Mức: VD
Hai lực vuông góc đồng quy có độ lớn lần lượt 5 $N$ và $5\,\text{N}$. Tính độ lớn hợp lực (kết quả theo đơn vị N).
\choiceTF
{\True Giá trị cần tìm là $7.07\,\text{N}$.}
{Giá trị cần tìm là $8.07\,\text{N}$.}
{\True Phải xác định đầy đủ các lực hoặc đại lượng liên quan trước khi tính.}
{Có thể bỏ qua điều kiện áp dụng của công thức.}
\loigiai{
\begin{itemchoice}
\item (Đúng) Giá trị cần tìm là $7.07\,\text{N}$. (Vì): kết quả $7.07\,\text{N}$. b) Sai. c) Đúng. d) Sai vì phải kiểm tra điều kiện áp dụng
\item (Sai) Giá trị cần tìm là $8.07\,\text{N}$.
\item (Đúng) Phải xác định đầy đủ các lực hoặc đại lượng liên quan trước khi tính.
\item (Sai) Có thể bỏ qua điều kiện áp dụng của công thức.
\end{itemchoice}
}
\end{ex}

% ===== Câu 88 =====
% ID: AIQ_L10C3_FULL_0096
% Mức: VD
\begin{ex}
% ID: AIQ_L10C3_FULL_0096
% Mức: VD
Hai lực vuông góc đồng quy có độ lớn lần lượt 6 $N$ và $6\,\text{N}$. Tính độ lớn hợp lực (kết quả theo đơn vị N).
\choiceTF
{\True Giá trị cần tìm là $8.49\,\text{N}$.}
{Giá trị cần tìm là $9.49\,\text{N}$.}
{\True Phải xác định đầy đủ các lực hoặc đại lượng liên quan trước khi tính.}
{Có thể bỏ qua điều kiện áp dụng của công thức.}
\loigiai{
\begin{itemchoice}
\item (Đúng) Giá trị cần tìm là $8.49\,\text{N}$. (Vì): kết quả $8.49\,\text{N}$. b) Sai. c) Đúng. d) Sai vì phải kiểm tra điều kiện áp dụng
\item (Sai) Giá trị cần tìm là $9.49\,\text{N}$.
\item (Đúng) Phải xác định đầy đủ các lực hoặc đại lượng liên quan trước khi tính.
\item (Sai) Có thể bỏ qua điều kiện áp dụng của công thức.
\end{itemchoice}
}
\end{ex}

% ===== Câu 89 =====
% ID: AIQ_L10C3_FULL_0097
% Mức: TH
\begin{ex}
% ID: AIQ_L10C3_FULL_0097
% Mức: TH
Hai lực vuông góc đồng quy có độ lớn lần lượt 3 $N$ và $4\,\text{N}$. Tính độ lớn hợp lực (kết quả theo đơn vị N).
\shortans{5}
\loigiai{
Thay số vào hệ thức của bài học thu được $5\,\text{N}$.
}
\end{ex}

% ===== Câu 90 =====
% ID: AIQ_L10C3_FULL_0098
% Mức: TH
\begin{ex}
% ID: AIQ_L10C3_FULL_0098
% Mức: TH
Hai lực vuông góc đồng quy có độ lớn lần lượt 4 $N$ và $5\,\text{N}$. Tính độ lớn hợp lực (kết quả theo đơn vị N).
\shortans{6.4}
\loigiai{
Thay số vào hệ thức của bài học thu được $6.4\,\text{N}$.
}
\end{ex}

% ===== Câu 91 =====
% ID: AIQ_L10C3_FULL_0099
% Mức: TH
\begin{ex}
% ID: AIQ_L10C3_FULL_0099
% Mức: TH
Hai lực vuông góc đồng quy có độ lớn lần lượt 5 $N$ và $6\,\text{N}$. Tính độ lớn hợp lực (kết quả theo đơn vị N).
\shortans{7.81}
\loigiai{
Thay số vào hệ thức của bài học thu được $7.81\,\text{N}$.
}
\end{ex}

% ===== Câu 92 =====
% ID: AIQ_L10C3_FULL_0100
% Mức: VD
\begin{ex}
% ID: AIQ_L10C3_FULL_0100
% Mức: VD
Hai lực vuông góc đồng quy có độ lớn lần lượt 6 $N$ và $4\,\text{N}$. Tính độ lớn hợp lực (kết quả theo đơn vị N).
\shortans{7.21}
\loigiai{
Thay số vào hệ thức của bài học thu được $7.21\,\text{N}$.
}
\end{ex}

% ===== Câu 93 =====
% ID: AIQ_L10C3_FULL_0101
% Mức: VD
\begin{ex}
% ID: AIQ_L10C3_FULL_0101
% Mức: VD
Hai lực vuông góc đồng quy có độ lớn lần lượt 3 $N$ và $5\,\text{N}$. Tính độ lớn hợp lực (kết quả theo đơn vị N).
\shortans{5.83}
\loigiai{
Thay số vào hệ thức của bài học thu được $5.83\,\text{N}$.
}
\end{ex}

% ===== Câu 94 =====
% ID: AIQ_L10C3_FULL_0102
% Mức: VD
\begin{ex}
% ID: AIQ_L10C3_FULL_0102
% Mức: VD
Hai lực vuông góc đồng quy có độ lớn lần lượt 4 $N$ và $6\,\text{N}$. Tính độ lớn hợp lực (kết quả theo đơn vị N).
\shortans{7.21}
\loigiai{
Thay số vào hệ thức của bài học thu được $7.21\,\text{N}$.
}
\end{ex}

% ===== Câu 95 =====
% ID: AIQ_L10C3_FULL_0103
% Mức: VD
\begin{bt}
% ID: AIQ_L10C3_FULL_0103
% Mức: VD
Hai lực vuông góc đồng quy có độ lớn lần lượt 5 $N$ và $4\,\text{N}$. Tính độ lớn hợp lực (kết quả theo đơn vị N).
\loigiai{
Xác định các đại lượng đã cho, chọn hệ thức phù hợp và thay số. Kết quả: $6.4\,\text{N}$.
}
\end{bt}

% ===== Câu 96 =====
% ID: AIQ_L10C3_FULL_0104
% Mức: VD
\begin{bt}
% ID: AIQ_L10C3_FULL_0104
% Mức: VD
Hai lực vuông góc đồng quy có độ lớn lần lượt 6 $N$ và $5\,\text{N}$. Tính độ lớn hợp lực (kết quả theo đơn vị N).
\loigiai{
Xác định các đại lượng đã cho, chọn hệ thức phù hợp và thay số. Kết quả: $7.81\,\text{N}$.
}
\end{bt}

% ===== Câu 97 =====
% ID: AIQ_L10C3_FULL_0105
% Mức: VD
\begin{bt}
% ID: AIQ_L10C3_FULL_0105
% Mức: VD
Hai lực vuông góc đồng quy có độ lớn lần lượt 3 $N$ và $6\,\text{N}$. Tính độ lớn hợp lực (kết quả theo đơn vị N).
\loigiai{
Xác định các đại lượng đã cho, chọn hệ thức phù hợp và thay số. Kết quả: $6.71\,\text{N}$.
}
\end{bt}

% ===== Câu 98 =====
% ID: AIQ_L10C3_FULL_0106
% Mức: VDC
\begin{bt}
% ID: AIQ_L10C3_FULL_0106
% Mức: VDC
Hai lực vuông góc đồng quy có độ lớn lần lượt 4 $N$ và $4\,\text{N}$. Tính độ lớn hợp lực (kết quả theo đơn vị N).
\loigiai{
Xác định các đại lượng đã cho, chọn hệ thức phù hợp và thay số. Kết quả: $5.66\,\text{N}$.
}
\end{bt}

% ===== Câu 99 =====
% ID: AIQ_L10C3_FULL_0107
% Mức: VDC
\begin{bt}
% ID: AIQ_L10C3_FULL_0107
% Mức: VDC
Hai lực vuông góc đồng quy có độ lớn lần lượt 5 $N$ và $5\,\text{N}$. Tính độ lớn hợp lực (kết quả theo đơn vị N).
\loigiai{
Xác định các đại lượng đã cho, chọn hệ thức phù hợp và thay số. Kết quả: $7.07\,\text{N}$.
}
\end{bt}

% ===== Câu 100 =====
% ID: AIQ_L10C3_FULL_0108
% Mức: VDC
\begin{bt}
% ID: AIQ_L10C3_FULL_0108
% Mức: VDC
Hai lực vuông góc đồng quy có độ lớn lần lượt 6 $N$ và $6\,\text{N}$. Tính độ lớn hợp lực (kết quả theo đơn vị N).
\loigiai{
Xác định các đại lượng đã cho, chọn hệ thức phù hợp và thay số. Kết quả: $8.49\,\text{N}$.
}
\end{bt}

% ===== Câu 101 =====
% ID: AIQ_L10C3_FULL_0109
% Mức: NB
\dangbt{Điều kiện cân bằng của chất điểm}
\begin{ex}
% ID: AIQ_L10C3_FULL_0109
% Mức: NB
Hai lực vuông góc đồng quy có độ lớn lần lượt 3 $N$ và $4\,\text{N}$. Tính độ lớn hợp lực (kết quả theo đơn vị N).
\choice
{\True 5}
{6}
{4}
{10}
\loigiai{
Áp dụng hệ thức phù hợp của bài học, thay số được $5\,\text{N}$.
}
\end{ex}

% ===== Câu 102 =====
% ID: AIQ_L10C3_FULL_0110
% Mức: NB
\begin{ex}
% ID: AIQ_L10C3_FULL_0110
% Mức: NB
Hai lực vuông góc đồng quy có độ lớn lần lượt 4 $N$ và $5\,\text{N}$. Tính độ lớn hợp lực (kết quả theo đơn vị N).
\choice
{7.4}
{\True 6.4}
{5.4}
{12.8}
\loigiai{
Áp dụng hệ thức phù hợp của bài học, thay số được $6.4\,\text{N}$.
}
\end{ex}

% ===== Câu 103 =====
% ID: AIQ_L10C3_FULL_0111
% Mức: NB
\begin{ex}
% ID: AIQ_L10C3_FULL_0111
% Mức: NB
Hai lực vuông góc đồng quy có độ lớn lần lượt 5 $N$ và $6\,\text{N}$. Tính độ lớn hợp lực (kết quả theo đơn vị N).
\choice
{6.81}
{8.81}
{\True 7.81}
{15.62}
\loigiai{
Áp dụng hệ thức phù hợp của bài học, thay số được $7.81\,\text{N}$.
}
\end{ex}

% ===== Câu 104 =====
% ID: AIQ_L10C3_FULL_0112
% Mức: NB
\begin{ex}
% ID: AIQ_L10C3_FULL_0112
% Mức: NB
Hai lực vuông góc đồng quy có độ lớn lần lượt 6 $N$ và $4\,\text{N}$. Tính độ lớn hợp lực (kết quả theo đơn vị N).
\choice
{14.42}
{8.21}
{6.21}
{\True 7.21}
\loigiai{
Áp dụng hệ thức phù hợp của bài học, thay số được $7.21\,\text{N}$.
}
\end{ex}

% ===== Câu 105 =====
% ID: AIQ_L10C3_FULL_0113
% Mức: TH
\begin{ex}
% ID: AIQ_L10C3_FULL_0113
% Mức: TH
Hai lực vuông góc đồng quy có độ lớn lần lượt 3 $N$ và $5\,\text{N}$. Tính độ lớn hợp lực (kết quả theo đơn vị N).
\choice
{\True 5.83}
{6.83}
{4.83}
{11.66}
\loigiai{
Áp dụng hệ thức phù hợp của bài học, thay số được $5.83\,\text{N}$.
}
\end{ex}

% ===== Câu 106 =====
% ID: AIQ_L10C3_FULL_0114
% Mức: TH
\begin{ex}
% ID: AIQ_L10C3_FULL_0114
% Mức: TH
Hai lực vuông góc đồng quy có độ lớn lần lượt 4 $N$ và $6\,\text{N}$. Tính độ lớn hợp lực (kết quả theo đơn vị N).
\choice
{8.21}
{\True 7.21}
{6.21}
{14.42}
\loigiai{
Áp dụng hệ thức phù hợp của bài học, thay số được $7.21\,\text{N}$.
}
\end{ex}

% ===== Câu 107 =====
% ID: AIQ_L10C3_FULL_0115
% Mức: TH
\begin{ex}
% ID: AIQ_L10C3_FULL_0115
% Mức: TH
Hai lực vuông góc đồng quy có độ lớn lần lượt 5 $N$ và $4\,\text{N}$. Tính độ lớn hợp lực (kết quả theo đơn vị N).
\choice
{5.4}
{7.4}
{\True 6.4}
{12.8}
\loigiai{
Áp dụng hệ thức phù hợp của bài học, thay số được $6.4\,\text{N}$.
}
\end{ex}

% ===== Câu 108 =====
% ID: AIQ_L10C3_FULL_0116
% Mức: TH
\begin{ex}
% ID: AIQ_L10C3_FULL_0116
% Mức: TH
Hai lực vuông góc đồng quy có độ lớn lần lượt 6 $N$ và $5\,\text{N}$. Tính độ lớn hợp lực (kết quả theo đơn vị N).
\choice
{15.62}
{8.81}
{6.81}
{\True 7.81}
\loigiai{
Áp dụng hệ thức phù hợp của bài học, thay số được $7.81\,\text{N}$.
}
\end{ex}

% ===== Câu 109 =====
% ID: AIQ_L10C3_FULL_0117
% Mức: VD
\begin{ex}
% ID: AIQ_L10C3_FULL_0117
% Mức: VD
Hai lực vuông góc đồng quy có độ lớn lần lượt 3 $N$ và $6\,\text{N}$. Tính độ lớn hợp lực (kết quả theo đơn vị N).
\choice
{\True 6.71}
{7.71}
{5.71}
{13.42}
\loigiai{
Áp dụng hệ thức phù hợp của bài học, thay số được $6.71\,\text{N}$.
}
\end{ex}

% ===== Câu 110 =====
% ID: AIQ_L10C3_FULL_0118
% Mức: VD
\begin{ex}
% ID: AIQ_L10C3_FULL_0118
% Mức: VD
Hai lực vuông góc đồng quy có độ lớn lần lượt 4 $N$ và $4\,\text{N}$. Tính độ lớn hợp lực (kết quả theo đơn vị N).
\choice
{6.66}
{\True 5.66}
{4.66}
{11.32}
\loigiai{
Áp dụng hệ thức phù hợp của bài học, thay số được $5.66\,\text{N}$.
}
\end{ex}

% ===== Câu 111 =====
% ID: AIQ_L10C3_FULL_0119
% Mức: TH
\begin{ex}
% ID: AIQ_L10C3_FULL_0119
% Mức: TH
Hai lực vuông góc đồng quy có độ lớn lần lượt 5 $N$ và $5\,\text{N}$. Tính độ lớn hợp lực (kết quả theo đơn vị N).
\choiceTF
{\True Giá trị cần tìm là $7.07\,\text{N}$.}
{Giá trị cần tìm là $8.07\,\text{N}$.}
{\True Phải xác định đầy đủ các lực hoặc đại lượng liên quan trước khi tính.}
{Có thể bỏ qua điều kiện áp dụng của công thức.}
\loigiai{
\begin{itemchoice}
\item (Đúng) Giá trị cần tìm là $7.07\,\text{N}$. (Vì): kết quả $7.07\,\text{N}$. b) Sai. c) Đúng. d) Sai vì phải kiểm tra điều kiện áp dụng
\item (Sai) Giá trị cần tìm là $8.07\,\text{N}$.
\item (Đúng) Phải xác định đầy đủ các lực hoặc đại lượng liên quan trước khi tính.
\item (Sai) Có thể bỏ qua điều kiện áp dụng của công thức.
\end{itemchoice}
}
\end{ex}

% ===== Câu 112 =====
% ID: AIQ_L10C3_FULL_0120
% Mức: TH
\begin{ex}
% ID: AIQ_L10C3_FULL_0120
% Mức: TH
Hai lực vuông góc đồng quy có độ lớn lần lượt 6 $N$ và $6\,\text{N}$. Tính độ lớn hợp lực (kết quả theo đơn vị N).
\choiceTF
{\True Giá trị cần tìm là $8.49\,\text{N}$.}
{Giá trị cần tìm là $9.49\,\text{N}$.}
{\True Phải xác định đầy đủ các lực hoặc đại lượng liên quan trước khi tính.}
{Có thể bỏ qua điều kiện áp dụng của công thức.}
\loigiai{
\begin{itemchoice}
\item (Đúng) Giá trị cần tìm là $8.49\,\text{N}$. (Vì): kết quả $8.49\,\text{N}$. b) Sai. c) Đúng. d) Sai vì phải kiểm tra điều kiện áp dụng
\item (Sai) Giá trị cần tìm là $9.49\,\text{N}$.
\item (Đúng) Phải xác định đầy đủ các lực hoặc đại lượng liên quan trước khi tính.
\item (Sai) Có thể bỏ qua điều kiện áp dụng của công thức.
\end{itemchoice}
}
\end{ex}

% ===== Câu 113 =====
% ID: AIQ_L10C3_FULL_0121
% Mức: VD
\begin{ex}
% ID: AIQ_L10C3_FULL_0121
% Mức: VD
Hai lực vuông góc đồng quy có độ lớn lần lượt 3 $N$ và $4\,\text{N}$. Tính độ lớn hợp lực (kết quả theo đơn vị N).
\choiceTF
{\True Giá trị cần tìm là $5\,\text{N}$.}
{Giá trị cần tìm là $6\,\text{N}$.}
{\True Phải xác định đầy đủ các lực hoặc đại lượng liên quan trước khi tính.}
{Có thể bỏ qua điều kiện áp dụng của công thức.}
\loigiai{
\begin{itemchoice}
\item (Đúng) Giá trị cần tìm là $5\,\text{N}$. (Vì): kết quả $5\,\text{N}$. b) Sai. c) Đúng. d) Sai vì phải kiểm tra điều kiện áp dụng
\item (Sai) Giá trị cần tìm là $6\,\text{N}$.
\item (Đúng) Phải xác định đầy đủ các lực hoặc đại lượng liên quan trước khi tính.
\item (Sai) Có thể bỏ qua điều kiện áp dụng của công thức.
\end{itemchoice}
}
\end{ex}

% ===== Câu 114 =====
% ID: AIQ_L10C3_FULL_0122
% Mức: VD
\begin{ex}
% ID: AIQ_L10C3_FULL_0122
% Mức: VD
Hai lực vuông góc đồng quy có độ lớn lần lượt 4 $N$ và $5\,\text{N}$. Tính độ lớn hợp lực (kết quả theo đơn vị N).
\choiceTF
{\True Giá trị cần tìm là $6.4\,\text{N}$.}
{Giá trị cần tìm là $7.4\,\text{N}$.}
{\True Phải xác định đầy đủ các lực hoặc đại lượng liên quan trước khi tính.}
{Có thể bỏ qua điều kiện áp dụng của công thức.}
\loigiai{
\begin{itemchoice}
\item (Đúng) Giá trị cần tìm là $6.4\,\text{N}$. (Vì): kết quả $6.4\,\text{N}$. b) Sai. c) Đúng. d) Sai vì phải kiểm tra điều kiện áp dụng
\item (Sai) Giá trị cần tìm là $7.4\,\text{N}$.
\item (Đúng) Phải xác định đầy đủ các lực hoặc đại lượng liên quan trước khi tính.
\item (Sai) Có thể bỏ qua điều kiện áp dụng của công thức.
\end{itemchoice}
}
\end{ex}

% ===== Câu 115 =====
% ID: AIQ_L10C3_FULL_0123
% Mức: VD
\begin{ex}
% ID: AIQ_L10C3_FULL_0123
% Mức: VD
Hai lực vuông góc đồng quy có độ lớn lần lượt 5 $N$ và $6\,\text{N}$. Tính độ lớn hợp lực (kết quả theo đơn vị N).
\choiceTF
{\True Giá trị cần tìm là $7.81\,\text{N}$.}
{Giá trị cần tìm là $8.81\,\text{N}$.}
{\True Phải xác định đầy đủ các lực hoặc đại lượng liên quan trước khi tính.}
{Có thể bỏ qua điều kiện áp dụng của công thức.}
\loigiai{
\begin{itemchoice}
\item (Đúng) Giá trị cần tìm là $7.81\,\text{N}$. (Vì): kết quả $7.81\,\text{N}$. b) Sai. c) Đúng. d) Sai vì phải kiểm tra điều kiện áp dụng
\item (Sai) Giá trị cần tìm là $8.81\,\text{N}$.
\item (Đúng) Phải xác định đầy đủ các lực hoặc đại lượng liên quan trước khi tính.
\item (Sai) Có thể bỏ qua điều kiện áp dụng của công thức.
\end{itemchoice}
}
\end{ex}

% ===== Câu 116 =====
% ID: AIQ_L10C3_FULL_0124
% Mức: TH
\begin{ex}
% ID: AIQ_L10C3_FULL_0124
% Mức: TH
Hai lực vuông góc đồng quy có độ lớn lần lượt 6 $N$ và $4\,\text{N}$. Tính độ lớn hợp lực (kết quả theo đơn vị N).
\shortans{7.21}
\loigiai{
Thay số vào hệ thức của bài học thu được $7.21\,\text{N}$.
}
\end{ex}

% ===== Câu 117 =====
% ID: AIQ_L10C3_FULL_0125
% Mức: TH
\begin{ex}
% ID: AIQ_L10C3_FULL_0125
% Mức: TH
Hai lực vuông góc đồng quy có độ lớn lần lượt 3 $N$ và $5\,\text{N}$. Tính độ lớn hợp lực (kết quả theo đơn vị N).
\shortans{5.83}
\loigiai{
Thay số vào hệ thức của bài học thu được $5.83\,\text{N}$.
}
\end{ex}

% ===== Câu 118 =====
% ID: AIQ_L10C3_FULL_0126
% Mức: TH
\begin{ex}
% ID: AIQ_L10C3_FULL_0126
% Mức: TH
Hai lực vuông góc đồng quy có độ lớn lần lượt 4 $N$ và $6\,\text{N}$. Tính độ lớn hợp lực (kết quả theo đơn vị N).
\shortans{7.21}
\loigiai{
Thay số vào hệ thức của bài học thu được $7.21\,\text{N}$.
}
\end{ex}

% ===== Câu 119 =====
% ID: AIQ_L10C3_FULL_0127
% Mức: VD
\begin{ex}
% ID: AIQ_L10C3_FULL_0127
% Mức: VD
Hai lực vuông góc đồng quy có độ lớn lần lượt 5 $N$ và $4\,\text{N}$. Tính độ lớn hợp lực (kết quả theo đơn vị N).
\shortans{6.4}
\loigiai{
Thay số vào hệ thức của bài học thu được $6.4\,\text{N}$.
}
\end{ex}

% ===== Câu 120 =====
% ID: AIQ_L10C3_FULL_0128
% Mức: VD
\begin{ex}
% ID: AIQ_L10C3_FULL_0128
% Mức: VD
Hai lực vuông góc đồng quy có độ lớn lần lượt 6 $N$ và $5\,\text{N}$. Tính độ lớn hợp lực (kết quả theo đơn vị N).
\shortans{7.81}
\loigiai{
Thay số vào hệ thức của bài học thu được $7.81\,\text{N}$.
}
\end{ex}

% ===== Câu 121 =====
% ID: AIQ_L10C3_FULL_0129
% Mức: VD
\begin{ex}
% ID: AIQ_L10C3_FULL_0129
% Mức: VD
Hai lực vuông góc đồng quy có độ lớn lần lượt 3 $N$ và $6\,\text{N}$. Tính độ lớn hợp lực (kết quả theo đơn vị N).
\shortans{6.71}
\loigiai{
Thay số vào hệ thức của bài học thu được $6.71\,\text{N}$.
}
\end{ex}

% ===== Câu 122 =====
% ID: AIQ_L10C3_FULL_0130
% Mức: VD
\begin{bt}
% ID: AIQ_L10C3_FULL_0130
% Mức: VD
Hai lực vuông góc đồng quy có độ lớn lần lượt 4 $N$ và $4\,\text{N}$. Tính độ lớn hợp lực (kết quả theo đơn vị N).
\loigiai{
Xác định các đại lượng đã cho, chọn hệ thức phù hợp và thay số. Kết quả: $5.66\,\text{N}$.
}
\end{bt}

% ===== Câu 123 =====
% ID: AIQ_L10C3_FULL_0131
% Mức: VD
\begin{bt}
% ID: AIQ_L10C3_FULL_0131
% Mức: VD
Hai lực vuông góc đồng quy có độ lớn lần lượt 5 $N$ và $5\,\text{N}$. Tính độ lớn hợp lực (kết quả theo đơn vị N).
\loigiai{
Xác định các đại lượng đã cho, chọn hệ thức phù hợp và thay số. Kết quả: $7.07\,\text{N}$.
}
\end{bt}

% ===== Câu 124 =====
% ID: AIQ_L10C3_FULL_0132
% Mức: VD
\begin{bt}
% ID: AIQ_L10C3_FULL_0132
% Mức: VD
Hai lực vuông góc đồng quy có độ lớn lần lượt 6 $N$ và $6\,\text{N}$. Tính độ lớn hợp lực (kết quả theo đơn vị N).
\loigiai{
Xác định các đại lượng đã cho, chọn hệ thức phù hợp và thay số. Kết quả: $8.49\,\text{N}$.
}
\end{bt}

% ===== Câu 125 =====
% ID: AIQ_L10C3_FULL_0133
% Mức: VDC
\begin{bt}
% ID: AIQ_L10C3_FULL_0133
% Mức: VDC
Hai lực vuông góc đồng quy có độ lớn lần lượt 3 $N$ và $4\,\text{N}$. Tính độ lớn hợp lực (kết quả theo đơn vị N).
\loigiai{
Xác định các đại lượng đã cho, chọn hệ thức phù hợp và thay số. Kết quả: $5\,\text{N}$.
}
\end{bt}

% ===== Câu 126 =====
% ID: AIQ_L10C3_FULL_0134
% Mức: VDC
\begin{bt}
% ID: AIQ_L10C3_FULL_0134
% Mức: VDC
Hai lực vuông góc đồng quy có độ lớn lần lượt 4 $N$ và $5\,\text{N}$. Tính độ lớn hợp lực (kết quả theo đơn vị N).
\loigiai{
Xác định các đại lượng đã cho, chọn hệ thức phù hợp và thay số. Kết quả: $6.4\,\text{N}$.
}
\end{bt}

% ===== Câu 127 =====
% ID: AIQ_L10C3_FULL_0135
% Mức: VDC
\begin{bt}
% ID: AIQ_L10C3_FULL_0135
% Mức: VDC
Hai lực vuông góc đồng quy có độ lớn lần lượt 5 $N$ và $6\,\text{N}$. Tính độ lớn hợp lực (kết quả theo đơn vị N).
\loigiai{
Xác định các đại lượng đã cho, chọn hệ thức phù hợp và thay số. Kết quả: $7.81\,\text{N}$.
}
\end{bt}

% ===== Câu 128 =====
% ID: AIQ_L10C3_FULL_0136
% Mức: NB
\dangbt{Bài toán tổng hợp và phân tích lực gắn với thực tiễn}
\begin{ex}
% ID: AIQ_L10C3_FULL_0136
% Mức: NB
Hai lực vuông góc đồng quy có độ lớn lần lượt 6 $N$ và $4\,\text{N}$. Tính độ lớn hợp lực (kết quả theo đơn vị N).
\choice
{\True 7.21}
{8.21}
{6.21}
{14.42}
\loigiai{
Áp dụng hệ thức phù hợp của bài học, thay số được $7.21\,\text{N}$.
}
\end{ex}

% ===== Câu 129 =====
% ID: AIQ_L10C3_FULL_0137
% Mức: NB
\begin{ex}
% ID: AIQ_L10C3_FULL_0137
% Mức: NB
Hai lực vuông góc đồng quy có độ lớn lần lượt 3 $N$ và $5\,\text{N}$. Tính độ lớn hợp lực (kết quả theo đơn vị N).
\choice
{6.83}
{\True 5.83}
{4.83}
{11.66}
\loigiai{
Áp dụng hệ thức phù hợp của bài học, thay số được $5.83\,\text{N}$.
}
\end{ex}

% ===== Câu 130 =====
% ID: AIQ_L10C3_FULL_0138
% Mức: NB
\begin{ex}
% ID: AIQ_L10C3_FULL_0138
% Mức: NB
Hai lực vuông góc đồng quy có độ lớn lần lượt 4 $N$ và $6\,\text{N}$. Tính độ lớn hợp lực (kết quả theo đơn vị N).
\choice
{6.21}
{8.21}
{\True 7.21}
{14.42}
\loigiai{
Áp dụng hệ thức phù hợp của bài học, thay số được $7.21\,\text{N}$.
}
\end{ex}

% ===== Câu 131 =====
% ID: AIQ_L10C3_FULL_0139
% Mức: NB
\begin{ex}
% ID: AIQ_L10C3_FULL_0139
% Mức: NB
Hai lực vuông góc đồng quy có độ lớn lần lượt 5 $N$ và $4\,\text{N}$. Tính độ lớn hợp lực (kết quả theo đơn vị N).
\choice
{12.8}
{7.4}
{5.4}
{\True 6.4}
\loigiai{
Áp dụng hệ thức phù hợp của bài học, thay số được $6.4\,\text{N}$.
}
\end{ex}

% ===== Câu 132 =====
% ID: AIQ_L10C3_FULL_0140
% Mức: TH
\begin{ex}
% ID: AIQ_L10C3_FULL_0140
% Mức: TH
Hai lực vuông góc đồng quy có độ lớn lần lượt 6 $N$ và $5\,\text{N}$. Tính độ lớn hợp lực (kết quả theo đơn vị N).
\choice
{\True 7.81}
{8.81}
{6.81}
{15.62}
\loigiai{
Áp dụng hệ thức phù hợp của bài học, thay số được $7.81\,\text{N}$.
}
\end{ex}

% ===== Câu 133 =====
% ID: AIQ_L10C3_FULL_0141
% Mức: TH
\begin{ex}
% ID: AIQ_L10C3_FULL_0141
% Mức: TH
Hai lực vuông góc đồng quy có độ lớn lần lượt 3 $N$ và $6\,\text{N}$. Tính độ lớn hợp lực (kết quả theo đơn vị N).
\choice
{7.71}
{\True 6.71}
{5.71}
{13.42}
\loigiai{
Áp dụng hệ thức phù hợp của bài học, thay số được $6.71\,\text{N}$.
}
\end{ex}

% ===== Câu 134 =====
% ID: AIQ_L10C3_FULL_0142
% Mức: TH
\begin{ex}
% ID: AIQ_L10C3_FULL_0142
% Mức: TH
Hai lực vuông góc đồng quy có độ lớn lần lượt 4 $N$ và $4\,\text{N}$. Tính độ lớn hợp lực (kết quả theo đơn vị N).
\choice
{4.66}
{6.66}
{\True 5.66}
{11.32}
\loigiai{
Áp dụng hệ thức phù hợp của bài học, thay số được $5.66\,\text{N}$.
}
\end{ex}

% ===== Câu 135 =====
% ID: AIQ_L10C3_FULL_0143
% Mức: TH
\begin{ex}
% ID: AIQ_L10C3_FULL_0143
% Mức: TH
Hai lực vuông góc đồng quy có độ lớn lần lượt 5 $N$ và $5\,\text{N}$. Tính độ lớn hợp lực (kết quả theo đơn vị N).
\choice
{14.14}
{8.07}
{6.07}
{\True 7.07}
\loigiai{
Áp dụng hệ thức phù hợp của bài học, thay số được $7.07\,\text{N}$.
}
\end{ex}

% ===== Câu 136 =====
% ID: AIQ_L10C3_FULL_0144
% Mức: VD
\begin{ex}
% ID: AIQ_L10C3_FULL_0144
% Mức: VD
Hai lực vuông góc đồng quy có độ lớn lần lượt 6 $N$ và $6\,\text{N}$. Tính độ lớn hợp lực (kết quả theo đơn vị N).
\choice
{\True 8.49}
{9.49}
{7.49}
{16.98}
\loigiai{
Áp dụng hệ thức phù hợp của bài học, thay số được $8.49\,\text{N}$.
}
\end{ex}

% ===== Câu 137 =====
% ID: AIQ_L10C3_FULL_0145
% Mức: VD
\begin{ex}
% ID: AIQ_L10C3_FULL_0145
% Mức: VD
Hai lực vuông góc đồng quy có độ lớn lần lượt 3 $N$ và $4\,\text{N}$. Tính độ lớn hợp lực (kết quả theo đơn vị N).
\choice
{6}
{\True 5}
{4}
{10}
\loigiai{
Áp dụng hệ thức phù hợp của bài học, thay số được $5\,\text{N}$.
}
\end{ex}

% ===== Câu 138 =====
% ID: AIQ_L10C3_FULL_0146
% Mức: TH
\begin{ex}
% ID: AIQ_L10C3_FULL_0146
% Mức: TH
Hai lực vuông góc đồng quy có độ lớn lần lượt 4 $N$ và $5\,\text{N}$. Tính độ lớn hợp lực (kết quả theo đơn vị N).
\choiceTF
{\True Giá trị cần tìm là $6.4\,\text{N}$.}
{Giá trị cần tìm là $7.4\,\text{N}$.}
{\True Phải xác định đầy đủ các lực hoặc đại lượng liên quan trước khi tính.}
{Có thể bỏ qua điều kiện áp dụng của công thức.}
\loigiai{
\begin{itemchoice}
\item (Đúng) Giá trị cần tìm là $6.4\,\text{N}$. (Vì): kết quả $6.4\,\text{N}$. b) Sai. c) Đúng. d) Sai vì phải kiểm tra điều kiện áp dụng
\item (Sai) Giá trị cần tìm là $7.4\,\text{N}$.
\item (Đúng) Phải xác định đầy đủ các lực hoặc đại lượng liên quan trước khi tính.
\item (Sai) Có thể bỏ qua điều kiện áp dụng của công thức.
\end{itemchoice}
}
\end{ex}

% ===== Câu 139 =====
% ID: AIQ_L10C3_FULL_0147
% Mức: TH
\begin{ex}
% ID: AIQ_L10C3_FULL_0147
% Mức: TH
Hai lực vuông góc đồng quy có độ lớn lần lượt 5 $N$ và $6\,\text{N}$. Tính độ lớn hợp lực (kết quả theo đơn vị N).
\choiceTF
{\True Giá trị cần tìm là $7.81\,\text{N}$.}
{Giá trị cần tìm là $8.81\,\text{N}$.}
{\True Phải xác định đầy đủ các lực hoặc đại lượng liên quan trước khi tính.}
{Có thể bỏ qua điều kiện áp dụng của công thức.}
\loigiai{
\begin{itemchoice}
\item (Đúng) Giá trị cần tìm là $7.81\,\text{N}$. (Vì): kết quả $7.81\,\text{N}$. b) Sai. c) Đúng. d) Sai vì phải kiểm tra điều kiện áp dụng
\item (Sai) Giá trị cần tìm là $8.81\,\text{N}$.
\item (Đúng) Phải xác định đầy đủ các lực hoặc đại lượng liên quan trước khi tính.
\item (Sai) Có thể bỏ qua điều kiện áp dụng của công thức.
\end{itemchoice}
}
\end{ex}

% ===== Câu 140 =====
% ID: AIQ_L10C3_FULL_0148
% Mức: VD
\begin{ex}
% ID: AIQ_L10C3_FULL_0148
% Mức: VD
Hai lực vuông góc đồng quy có độ lớn lần lượt 6 $N$ và $4\,\text{N}$. Tính độ lớn hợp lực (kết quả theo đơn vị N).
\choiceTF
{\True Giá trị cần tìm là $7.21\,\text{N}$.}
{Giá trị cần tìm là $8.21\,\text{N}$.}
{\True Phải xác định đầy đủ các lực hoặc đại lượng liên quan trước khi tính.}
{Có thể bỏ qua điều kiện áp dụng của công thức.}
\loigiai{
\begin{itemchoice}
\item (Đúng) Giá trị cần tìm là $7.21\,\text{N}$. (Vì): kết quả $7.21\,\text{N}$. b) Sai. c) Đúng. d) Sai vì phải kiểm tra điều kiện áp dụng
\item (Sai) Giá trị cần tìm là $8.21\,\text{N}$.
\item (Đúng) Phải xác định đầy đủ các lực hoặc đại lượng liên quan trước khi tính.
\item (Sai) Có thể bỏ qua điều kiện áp dụng của công thức.
\end{itemchoice}
}
\end{ex}

% ===== Câu 141 =====
% ID: AIQ_L10C3_FULL_0149
% Mức: VD
\begin{ex}
% ID: AIQ_L10C3_FULL_0149
% Mức: VD
Hai lực vuông góc đồng quy có độ lớn lần lượt 3 $N$ và $5\,\text{N}$. Tính độ lớn hợp lực (kết quả theo đơn vị N).
\choiceTF
{\True Giá trị cần tìm là $5.83\,\text{N}$.}
{Giá trị cần tìm là $6.83\,\text{N}$.}
{\True Phải xác định đầy đủ các lực hoặc đại lượng liên quan trước khi tính.}
{Có thể bỏ qua điều kiện áp dụng của công thức.}
\loigiai{
\begin{itemchoice}
\item (Đúng) Giá trị cần tìm là $5.83\,\text{N}$. (Vì): kết quả $5.83\,\text{N}$. b) Sai. c) Đúng. d) Sai vì phải kiểm tra điều kiện áp dụng
\item (Sai) Giá trị cần tìm là $6.83\,\text{N}$.
\item (Đúng) Phải xác định đầy đủ các lực hoặc đại lượng liên quan trước khi tính.
\item (Sai) Có thể bỏ qua điều kiện áp dụng của công thức.
\end{itemchoice}
}
\end{ex}

% ===== Câu 142 =====
% ID: AIQ_L10C3_FULL_0150
% Mức: VD
\begin{ex}
% ID: AIQ_L10C3_FULL_0150
% Mức: VD
Hai lực vuông góc đồng quy có độ lớn lần lượt 4 $N$ và $6\,\text{N}$. Tính độ lớn hợp lực (kết quả theo đơn vị N).
\choiceTF
{\True Giá trị cần tìm là $7.21\,\text{N}$.}
{Giá trị cần tìm là $8.21\,\text{N}$.}
{\True Phải xác định đầy đủ các lực hoặc đại lượng liên quan trước khi tính.}
{Có thể bỏ qua điều kiện áp dụng của công thức.}
\loigiai{
\begin{itemchoice}
\item (Đúng) Giá trị cần tìm là $7.21\,\text{N}$. (Vì): kết quả $7.21\,\text{N}$. b) Sai. c) Đúng. d) Sai vì phải kiểm tra điều kiện áp dụng
\item (Sai) Giá trị cần tìm là $8.21\,\text{N}$.
\item (Đúng) Phải xác định đầy đủ các lực hoặc đại lượng liên quan trước khi tính.
\item (Sai) Có thể bỏ qua điều kiện áp dụng của công thức.
\end{itemchoice}
}
\end{ex}

% ===== Câu 143 =====
% ID: AIQ_L10C3_FULL_0151
% Mức: TH
\begin{ex}
% ID: AIQ_L10C3_FULL_0151
% Mức: TH
Hai lực vuông góc đồng quy có độ lớn lần lượt 5 $N$ và $4\,\text{N}$. Tính độ lớn hợp lực (kết quả theo đơn vị N).
\shortans{6.4}
\loigiai{
Thay số vào hệ thức của bài học thu được $6.4\,\text{N}$.
}
\end{ex}

% ===== Câu 144 =====
% ID: AIQ_L10C3_FULL_0152
% Mức: TH
\begin{ex}
% ID: AIQ_L10C3_FULL_0152
% Mức: TH
Hai lực vuông góc đồng quy có độ lớn lần lượt 6 $N$ và $5\,\text{N}$. Tính độ lớn hợp lực (kết quả theo đơn vị N).
\shortans{7.81}
\loigiai{
Thay số vào hệ thức của bài học thu được $7.81\,\text{N}$.
}
\end{ex}

% ===== Câu 145 =====
% ID: AIQ_L10C3_FULL_0153
% Mức: TH
\begin{ex}
% ID: AIQ_L10C3_FULL_0153
% Mức: TH
Hai lực vuông góc đồng quy có độ lớn lần lượt 3 $N$ và $6\,\text{N}$. Tính độ lớn hợp lực (kết quả theo đơn vị N).
\shortans{6.71}
\loigiai{
Thay số vào hệ thức của bài học thu được $6.71\,\text{N}$.
}
\end{ex}

% ===== Câu 146 =====
% ID: AIQ_L10C3_FULL_0154
% Mức: VD
\begin{ex}
% ID: AIQ_L10C3_FULL_0154
% Mức: VD
Hai lực vuông góc đồng quy có độ lớn lần lượt 4 $N$ và $4\,\text{N}$. Tính độ lớn hợp lực (kết quả theo đơn vị N).
\shortans{5.66}
\loigiai{
Thay số vào hệ thức của bài học thu được $5.66\,\text{N}$.
}
\end{ex}

% ===== Câu 147 =====
% ID: AIQ_L10C3_FULL_0155
% Mức: VD
\begin{ex}
% ID: AIQ_L10C3_FULL_0155
% Mức: VD
Hai lực vuông góc đồng quy có độ lớn lần lượt 5 $N$ và $5\,\text{N}$. Tính độ lớn hợp lực (kết quả theo đơn vị N).
\shortans{7.07}
\loigiai{
Thay số vào hệ thức của bài học thu được $7.07\,\text{N}$.
}
\end{ex}

% ===== Câu 148 =====
% ID: AIQ_L10C3_FULL_0156
% Mức: VD
\begin{ex}
% ID: AIQ_L10C3_FULL_0156
% Mức: VD
Hai lực vuông góc đồng quy có độ lớn lần lượt 6 $N$ và $6\,\text{N}$. Tính độ lớn hợp lực (kết quả theo đơn vị N).
\shortans{8.49}
\loigiai{
Thay số vào hệ thức của bài học thu được $8.49\,\text{N}$.
}
\end{ex}

% ===== Câu 149 =====
% ID: AIQ_L10C3_FULL_0157
% Mức: VD
\begin{bt}
% ID: AIQ_L10C3_FULL_0157
% Mức: VD
Hai lực vuông góc đồng quy có độ lớn lần lượt 3 $N$ và $4\,\text{N}$. Tính độ lớn hợp lực (kết quả theo đơn vị N).
\loigiai{
Xác định các đại lượng đã cho, chọn hệ thức phù hợp và thay số. Kết quả: $5\,\text{N}$.
}
\end{bt}

% ===== Câu 150 =====
% ID: AIQ_L10C3_FULL_0158
% Mức: VD
\begin{bt}
% ID: AIQ_L10C3_FULL_0158
% Mức: VD
Hai lực vuông góc đồng quy có độ lớn lần lượt 4 $N$ và $5\,\text{N}$. Tính độ lớn hợp lực (kết quả theo đơn vị N).
\loigiai{
Xác định các đại lượng đã cho, chọn hệ thức phù hợp và thay số. Kết quả: $6.4\,\text{N}$.
}
\end{bt}

% ===== Câu 151 =====
% ID: AIQ_L10C3_FULL_0159
% Mức: VD
\begin{bt}
% ID: AIQ_L10C3_FULL_0159
% Mức: VD
Hai lực vuông góc đồng quy có độ lớn lần lượt 5 $N$ và $6\,\text{N}$. Tính độ lớn hợp lực (kết quả theo đơn vị N).
\loigiai{
Xác định các đại lượng đã cho, chọn hệ thức phù hợp và thay số. Kết quả: $7.81\,\text{N}$.
}
\end{bt}

% ===== Câu 152 =====
% ID: AIQ_L10C3_FULL_0160
% Mức: VDC
\begin{bt}
% ID: AIQ_L10C3_FULL_0160
% Mức: VDC
Hai lực vuông góc đồng quy có độ lớn lần lượt 6 $N$ và $4\,\text{N}$. Tính độ lớn hợp lực (kết quả theo đơn vị N).
\loigiai{
Xác định các đại lượng đã cho, chọn hệ thức phù hợp và thay số. Kết quả: $7.21\,\text{N}$.
}
\end{bt}

% ===== Câu 153 =====
% ID: AIQ_L10C3_FULL_0161
% Mức: VDC
\begin{bt}
% ID: AIQ_L10C3_FULL_0161
% Mức: VDC
Hai lực vuông góc đồng quy có độ lớn lần lượt 3 $N$ và $5\,\text{N}$. Tính độ lớn hợp lực (kết quả theo đơn vị N).
\loigiai{
Xác định các đại lượng đã cho, chọn hệ thức phù hợp và thay số. Kết quả: $5.83\,\text{N}$.
}
\end{bt}

% ===== Câu 154 =====
% ID: AIQ_L10C3_FULL_0162
% Mức: VDC
\begin{bt}
% ID: AIQ_L10C3_FULL_0162
% Mức: VDC
Hai lực vuông góc đồng quy có độ lớn lần lượt 4 $N$ và $6\,\text{N}$. Tính độ lớn hợp lực (kết quả theo đơn vị N).
\loigiai{
Xác định các đại lượng đã cho, chọn hệ thức phù hợp và thay số. Kết quả: $7.21\,\text{N}$.
}
\end{bt}
